\documentclass[11pt]{amsart}

\usepackage[T1]{fontenc}
\usepackage{amsmath}
\usepackage{amssymb}
\usepackage{amsthm}
\usepackage{booktabs}
\usepackage{array}
\usepackage{xcolor}
\usepackage{fancyvrb}
\usepackage{url}
\usepackage{hyperref}

% ---------------------------------------------------------------------------
% Theorem environments
% ---------------------------------------------------------------------------
\newtheorem{theorem}{Theorem}[section]
\newtheorem{proposition}[theorem]{Proposition}
\newtheorem{lemma}[theorem]{Lemma}
\newtheorem{corollary}[theorem]{Corollary}
\theoremstyle{definition}
\newtheorem{definition}[theorem]{Definition}
\theoremstyle{remark}
\newtheorem{remark}[theorem]{Remark}

\numberwithin{equation}{section}
\setlength{\emergencystretch}{3em}

% ---------------------------------------------------------------------------
% Macros
% ---------------------------------------------------------------------------
\newcommand{\NN}{\mathbb{N}}
\newcommand{\ZZ}{\mathbb{Z}}
\newcommand{\QQ}{\mathbb{Q}}
\newcommand{\RR}{\mathbb{R}}
\newcommand{\CC}{\mathbb{C}}
\newcommand{\PP}{\mathbb{P}}
\newcommand{\Qbar}{\overline{\mathbb{Q}}}
\newcommand{\Ga}{\mathbb{G}_{\mathrm{a}}}
\newcommand{\Gm}{\mathbb{G}_{\mathrm{m}}}
\newcommand{\hdim}{\operatorname{hdim}}
\newcommand{\rk}{\operatorname{rk}}
\newcommand{\Lie}{\operatorname{Lie}}
\newcommand{\codim}{\operatorname{codim}}
\newcommand{\vp}{\varpi}
\newcommand{\lieAlg}{\operatorname{Lie}(H)}
\newcommand{\NR}{\textup{(NR)}}
\DeclareUrlCommand\lean{\urlstyle{tt}}

% ---------------------------------------------------------------------------
% Lean code blocks: fancyvrb, with the Unicode used in them mapped for pdflatex
% ---------------------------------------------------------------------------
\fvset{frame=single,fontsize=\footnotesize,xleftmargin=1em,framesep=2mm}
\DeclareUnicodeCharacter{2102}{\ensuremath{\mathbb{C}}}
\DeclareUnicodeCharacter{211A}{\ensuremath{\mathbb{Q}}}
\DeclareUnicodeCharacter{2115}{\ensuremath{\mathbb{N}}}
\DeclareUnicodeCharacter{2124}{\ensuremath{\mathbb{Z}}}
\DeclareUnicodeCharacter{211D}{\ensuremath{\mathbb{R}}}
\DeclareUnicodeCharacter{00D7}{\ensuremath{\times}}
\DeclareUnicodeCharacter{2200}{\ensuremath{\forall}}
\DeclareUnicodeCharacter{2203}{\ensuremath{\exists}}
\DeclareUnicodeCharacter{2192}{\ensuremath{\to}}
\DeclareUnicodeCharacter{2208}{\ensuremath{\in}}
\DeclareUnicodeCharacter{2209}{\ensuremath{\notin}}
\DeclareUnicodeCharacter{2227}{\ensuremath{\wedge}}
\DeclareUnicodeCharacter{2228}{\ensuremath{\vee}}
\DeclareUnicodeCharacter{2211}{\ensuremath{\sum}}
\DeclareUnicodeCharacter{2293}{\ensuremath{\sqcap}}
\DeclareUnicodeCharacter{2294}{\ensuremath{\sqcup}}
\DeclareUnicodeCharacter{2264}{\ensuremath{\le}}
\DeclareUnicodeCharacter{2260}{\ensuremath{\ne}}
\DeclareUnicodeCharacter{2265}{\ensuremath{\ge}}
\DeclareUnicodeCharacter{21A5}{\ensuremath{\uparrow}}
\DeclareUnicodeCharacter{2080}{\ensuremath{_0}}
\DeclareUnicodeCharacter{2081}{\ensuremath{_1}}
\DeclareUnicodeCharacter{2082}{\ensuremath{_2}}
\DeclareUnicodeCharacter{2083}{\ensuremath{_3}}
\DeclareUnicodeCharacter{03B4}{\ensuremath{\delta}}
\DeclareUnicodeCharacter{03B1}{\ensuremath{\alpha}}
\DeclareUnicodeCharacter{03BC}{\ensuremath{\mu}}
\DeclareUnicodeCharacter{03C9}{\ensuremath{\omega}}
\DeclareUnicodeCharacter{03BA}{\ensuremath{\kappa}}
\DeclareUnicodeCharacter{03C1}{\ensuremath{\rho}}
\DeclareUnicodeCharacter{03C0}{\ensuremath{\pi}}
\DeclareUnicodeCharacter{03C6}{\ensuremath{\varphi}}
\DeclareUnicodeCharacter{2032}{\ensuremath{'}}
\DeclareUnicodeCharacter{22A4}{\ensuremath{\top}}
\DeclareUnicodeCharacter{22A5}{\ensuremath{\bot}}

\begin{document}

\title[Waldschmidt's subgroup theorem with period term in Lean]{A Lean formalization of Waldschmidt's
subgroup theorem with period term for groups with elliptic factors}

\author{Shawn Krisman}

\subjclass[2020]{Primary 11J89; Secondary 11J81, 11J95, 11G05, 14L10, 68V20}
\keywords{transcendence, algebraic subgroup theorem, multiplicity estimates, elliptic functions,
Schneider's theorem, Schneider--Lang criterion, formal verification, Lean}

\date{\today}

\begin{abstract}
We formalize in Lean~4 with Mathlib Waldschmidt's algebraic subgroup theorem with period term
(Theorem~4.1 of his 1988 Durham paper) for the groups
$G=\Ga^{d_0}\times\Gm^{d_1}\times E_1\times\dots\times E_m$, where the $E_j$ are elliptic curves
with algebraic invariants, under the additional hypothesis $2N\ge d_1+2d_0+\kappa$ ($N=\dim V$,
$\kappa$ the rank of the periods in the data). The proof follows Waldschmidt: an auxiliary
function built with periods, which saves the factor $r^{-\kappa}$, combined with Philippon's
multiplicity estimate. Both are proved in Lean for elliptic factors; the multiplicity estimate
required several changes to the textbook route, which we describe. Along the way the development
gives what we believe are the first formal proofs of Schneider's theorems on elliptic functions
(periods, periods divided by $\pi$, and values of $\wp$ at algebraic points are transcendental) and of
the Schneider--Lang criterion for meromorphic functions, and the first Lean library for the
Weierstrass $\sigma$ and $\zeta$ functions, including Legendre's relation and the addition theorem.
The proof closure uses only Lean's three standard axioms. The theorem discharges, in a companion
certificate, the hypothesis behind the four exponentials statement at $t=c\vp/\pi$ for every
algebraic $c\ne0$. We also prove the theorem in restricted cases for groups
$\Ga^{d_0}\times\Gm^{d_1}\times G_2$ with $G_2$ the push-out of a product of universal vector
extensions $E_j^\natural$ along an algebraic $r\times m$ matrix: for $r\le1$ when $V$ meets the vector
fibre trivially, and for $d_0=0$, $r\le2$ when $V$ meets it in dimension at most one, for curves
satisfying two lifting properties that hold over $\Qbar$ and are proved in Lean for the CM curves
used. With these, all elliptic and vector-extension hypotheses of the companion certificate are
discharged, among them the one behind the four exponentials statement at $t=c\vp$. Finally, we
prove the theorem and its hyperplane consequence for $\Ga^{d_0}\times\Gm^{d_1}\times A^r$, where $A$ is the
Jacobian of $y^5=x(1-x)$, for toric-injective hyperplanes and data whose $A$-components are periods; this
needs a multiplicity estimate for abelian blocks given by an explicit projective model, whose polynomial
identities are checked by kernel-evaluated certificates. It discharges the hypothesis behind the
companion certificate's statements at the Beta value $B(1/5,1/5)$, for example that
$2^{B(1/5,1/5)/\pi}$ and $3^{B(1/5,1/5)/\pi}$ are not both algebraic. With the same two restrictions we also
prove the hyperplane consequence for a push-out of powers of a $\Ga$-extension of this Jacobian; in the companion
certificate this gives that $2^{B(1/5,1/5)}$ and $3^{B(1/5,1/5)}$ are not both algebraic.
\end{abstract}

\maketitle

% ===========================================================================
\section{Introduction}\label{sec:intro}
% ===========================================================================

This paper is a sequel to \cite{StrongSixPaper}, which formalized Roy's strong six exponentials
theorem in Lean. That development proved, as its main intermediate result, Waldschmidt's
transcendence theorem on linear groups $\Ga^{d_0}\times\Gm^{d_1}$ in the form of Roy's Theorem~1
\cite{Roy1992}, together with Philippon's multiplicity estimate on those groups. Here we leave the
linear world and treat groups with elliptic factors.

\subsection{The theorem}
Let $G$ be a connected commutative algebraic group over $\Qbar$ of dimension $d$, $T_G$ its tangent
space at the origin with its $\Qbar$-structure, $\exp_G\colon T_G\to G(\CC)$ the exponential map,
$\Omega(G)=\ker\exp_G$ and $\Lambda(G)=\exp_G^{-1}G(\Qbar)$. Write $G=\Ga^{d_0}\times\Gm^{d_1}\times G_2$.
Waldschmidt's Theorem~4.1 of \cite{W88} is the following.

\begin{theorem}[{Waldschmidt \cite[Thm.~4.1, pp.~385--386]{W88}}]\label{thm:W88}
Let $V\subset T_G$ be a subspace of dimension $N<d$, $W\subset V$ a subspace defined over $\Qbar$, and
$Y\subset V\cap\Lambda(G)$ a finitely generated subgroup; put $\kappa=\rk_\ZZ(Y\cap\Omega(G))$. Then
there is a connected algebraic subgroup $G'\ne G$ defined over $\Qbar$ with $\delta>\tau$ and
\begin{equation}\label{eq:W88}
(\lambda+\delta_1+2\delta_2)(d-N)\le(\delta-\tau)(d_1+2d_2-\kappa),
\end{equation}
where $\delta=\codim G'$, $\delta_0$ and $\delta_1$ are the codimensions of the projections of $G'$ to
$\Ga^{d_0}$ and $\Gm^{d_1}$, $\delta_2=\delta-\delta_0-\delta_1$, $\lambda=\rk_\ZZ Y/(Y\cap\exp_G^{-1}G')$ and
$\tau=\dim W/(W\cap T_{G'})$.
\end{theorem}

The \emph{period term} is $-\kappa$ on the right of \eqref{eq:W88}. For $\kappa=0$ the theorem is
the algebraic subgroup theorem in its usual form, a common source of the six exponentials theorem,
Baker's theorem and their elliptic analogues \cite{WaldschmidtLIL,DALAG}. The refinement by $\kappa$
matters when the data contain periods that do not span the Lie algebra of an algebraic subgroup,
for instance a period of $\Gm\times E$ lying on a line through a logarithm.

\subsection{What is formalized}
In Lean, the main theorem is
\[
\texttt{theorem LSG.W88.theorem41Abelian : LSG.W88.Theorem41Abelian},
\]
where the proposition \lean{Theorem41Abelian} is defined in \lean{LSG/W88/Statement.lean}. It is
Theorem~\ref{thm:W88} for
\[
G=\Ga^{d_0}\times\Gm^{d_1}\times E_1\times\dots\times E_m,\qquad E_j\colon y^2=4x^3-g_2(\Lambda_j)x-g_3(\Lambda_j),
\]
for arbitrary lattices $\Lambda_j$ with $g_2,g_3\in\Qbar$, under the extra hypothesis
\begin{equation}\label{eq:NR}
\NR\qquad 2N\ge d_1+2d_0+\kappa .
\end{equation}
Section~\ref{sec:statement} reproduces the statement. Hypothesis \NR{} is the range in which we prove
Waldschmidt's Proposition~6.1 (the auxiliary function) with multiplicities along $W$; it is the
condition of Lemma~U of \cite{MstarPaper}, which supplies the step of \cite[Prop.~6.1]{W88} that is
only indicated in print. \NR{} is a genuine restriction (Remark~\ref{rem:NR}). Correspondingly,
results of this paper must be cited as ``W88 Thm.~4.1 under \NR{} for
$\Ga^{d_0}\times\Gm^{d_1}\times E_1\times\dots\times E_m$''.

From the main theorem we derive, in Lean, its consequence of \cite[\S5d]{W88} (Theorem~1.1 with the
$\kappa$-refined hypothesis (5.4)) for $\Gm^{d_1}\times E$ and for $\Ga^{d_0}\times\Gm^{d_1}\times E^r$,
and Theorem~\ref{thm:W88} for $\Ga^{k_0}\times\Gm^{k_1}\times B$ with $B$ an arbitrary abelian
subvariety of a product of elliptic curves (Section~\ref{sec:5d}).

For groups whose factor $G_2$ is a push-out of $(E_1\times\dots\times E_m)^\natural$ along an algebraic
matrix we prove Theorem~\ref{thm:W88} under \NR{} and its \S5d consequence only in restricted cases,
listed in Theorem~\ref{thm:vec} (Section~\ref{sec:vecext}). Those results must be cited with their
scope, never as ``W88 Thm.~4.1 for vector extensions''.

For $G=\Ga^{d_0}\times\Gm^{d_1}\times A^r$ with $A$ the Jacobian of $y^5=x(1-x)$ we prove
Theorem~\ref{thm:W88} under \NR{} and its \S5d consequence for toric-injective hyperplanes and for subgroups
$Y$ whose $A$-components are (rational) periods (Theorem~\ref{thm:abvar}, Section~\ref{sec:abvar}). With the same
restrictions we prove the \S5d consequence for $\Ga^{d_0}\times\Gm^{d_1}\times Q$, $Q$ a push-out of powers of the
$\Ga$-extension $A_{\bar\sigma}^\natural$ of $A$ (Theorem~\ref{thm:abvarv}, Section~\ref{sec:abvarv}).

\subsection{Why it matters here: the M* applications}
The paper \cite{MstarPaper} proves four exponentials statements at CM periods by applying
Theorem~\ref{thm:W88} to explicit groups whose subgroups it lists: for a CM curve with period
$\Omega$ and $c\in\Qbar^\times$, if $\log\alpha,\log\beta$ are $\QQ$-linearly independent then
$\alpha^{c\Omega},\beta^{c\Omega}$ are not both algebraic, and neither are $\alpha^{c\Omega/\pi},
\beta^{c\Omega/\pi}$; for example $2^{\vp/\pi}$ and $3^{\vp/\pi}$ are not both algebraic, with
$\vp=\Gamma(1/4)^2/(2\sqrt{2\pi})$. A separate Lean certificate, \lean{lean_mstar}
\cite{LeanMstar}, checks these deductions with Waldschmidt's theorem for the relevant groups as
explicit hypotheses. The present work discharges those hypotheses whose groups are built from
tori, additive groups and powers of one elliptic curve: \lean{lean_mstar}'s \lean{Thm11Kappa} is our
\lean{LSG.W88.thm11Kappa} verbatim, and its \lean{HypStdCMAt} is \lean{LSG.W88.hypStdCMAt}. In
\lean{lean_mstar} this makes the four exponentials statement at $t=c\vp/\pi$, for every algebraic
$c\ne0$, and the paper's ``Theorem~M'' for $\Omega=c\vp$ unconditional (standard axioms only). The
statements at $t=c\vp$ and the two-curve results use groups with universal vector extensions
$E^\natural$. Their hypotheses fall within the restricted cases of Section~\ref{sec:vecext}, and are
discharged there as well; in \lean{lean_mstar} these results, too, no longer assume Waldschmidt's
theorem (some of them still use other named classical inputs). The statements at the Beta value
$B(1/5,1/5)$ (Theorem~A of \cite{MstarPaper} for the Jacobian of $y^5=x(1-x)$) use the groups of
Section~\ref{sec:abvar} and are discharged there. The headline statement at $t=B(1/5,1/5)$ itself uses the groups of
Section~\ref{sec:abvarv} and is discharged there, again on toric-injective hyperplanes.

\subsection{First formal proofs}
To reach the main theorem we had to build the classical theory of elliptic functions in Lean
beyond what Mathlib v4.35.0 provides (Mathlib has $\wp$, $\wp'$, the invariants $g_2,g_3$ and the
differential equation for a \lean{PeriodPair}). Section~\ref{sec:stage1} describes the library: the
Weierstrass $\sigma$ function as an infinite product, $\zeta$, the quasi-periods $\eta$,
quasi-periodicity of $\sigma$, growth of order two, Legendre's relation, the addition theorem and the
dictionary with Mathlib's algebraic group law on Weierstrass curves. On top of it we prove the
Schneider--Lang criterion for meromorphic functions of finite order in one variable and Schneider's
theorems \cite{Schneider1937}: if $g_2,g_3$ are algebraic, then every nonzero period $\omega$ and
$\omega/\pi$ are transcendental, and $\wp(\alpha)$ is transcendental for algebraic $\alpha\notin\Lambda$.

\emph{Prior formalizations.} We searched (on 30 October 2026: the web, the Mathlib v4.35.0 sources,
and the prior-art collections we keep; the web again, more briefly, on 6 November 2026) for formal proofs of transcendence results about elliptic
functions and found none in any proof assistant. The closest work we know: Karatarakis and
Wiedijk's Lean proof of the Gelfond--Schneider theorem \cite{KaratarakisWiedijk2026}; the library
\cite{DiazModulus}, which proves the Schneider--Lang criterion for $\CC^{d_0}\times(\CC^\times)^{d_1}$
(exponential functions only), Gelfond--Schneider and the qualitative form of Baker's theorem;
an Isabelle/HOL development of elliptic functions and modular forms ($\wp$, the addition theorem,
theta functions, the Jacobi triple product, Dedekind's $\eta$) \cite{EberlLiPaulson2026}, whose
abstract and keywords mention no transcendence results; and the Lean files of
\cite{ChCh26}, which prove irrationality-exponent bounds for CM periods with uniformization, theta
and CM inputs taken as hypotheses (Diophantine results, not transcendence; no $\sigma$, $\zeta$ or
quasi-periods). With this evidence we believe that the formalization contains the first formal
proofs of a transcendence theorem for elliptic functions and of the meromorphic Schneider--Lang
criterion, the first Lean development of $\sigma$ and $\zeta$, and the first formal multiplicity
estimate on a non-linear group. The search was not exhaustive, and we would be glad to learn of
earlier work.

\subsection{Organization}
Section~\ref{sec:statement} gives the formal statement and its review. Section~\ref{sec:arch}
describes the architecture. Sections~\ref{sec:stage1}--\ref{sec:5d} describe the three stages and
the bridges, with emphasis on the places where the formal proof had to depart from the literature.
Section~\ref{sec:vecext} treats the vector-extension cases, Section~\ref{sec:abvar} the Jacobian
of $y^5=x(1-x)$, Section~\ref{sec:abvarv} a vector extension of it, Section~\ref{sec:conseq} the explicit consequences,
Section~\ref{sec:stats} the artifact, and Section~\ref{sec:trust} what has and has not been checked.

% ===========================================================================
\section{The formal statement}\label{sec:statement}
% ===========================================================================

The statement uses no algebraic-group theory. Everything is expressed on the tangent space
$T_G=\CC^{d_0}\times\CC^{d_1}\times\CC^m$, with $\exp_G(a,b,z)=(a,e^b,(\wp_j(z_j),\wp_j'(z_j))_j)$. The
elliptic curves are Mathlib's: a \lean{PeriodPair} \lean{L} is a pair of $\RR$-independent complex
numbers, with lattice \lean{L.lattice}, Weierstrass function \lean{L.weierstrassP} and invariants
\texttt{L.g}$_2$, \texttt{L.g}$_3$.

\begin{Verbatim}[fontsize=\scriptsize]
abbrev Tan (d0 d1 m : ℕ) := (Fin d0 → ℂ) × (Fin d1 → ℂ) × (Fin m → ℂ)

def Lam (Ls : Fin m → PeriodPair) : Set (Tan d0 d1 m) :=
  {p | AlgVec p.1 ∧ (∀ i, IsAlgebraic ℚ (Complex.exp (p.2.1 i))) ∧
    ∀ j, p.2.2 j ∈ (Ls j).lattice ∨
      IsAlgebraic ℚ ((Ls j).weierstrassP (p.2.2 j))}

-- carrier of the AddSubgroup `Per Ls`:
  {p | p.1 = 0 ∧ (∀ i, ∃ k : ℤ, p.2.1 i = 2 * Real.pi * I * k) ∧
    ∀ j, p.2.2 j ∈ (Ls j).lattice}

def IsAbelianTangent (Ls : Fin m → PeriodPair)
    (t : Submodule ℂ (Fin m → ℂ)) : Prop :=
  (t.restrictScalars ℝ : Submodule ℝ (Fin m → ℂ)) =
    Submodule.span ℝ {z | z ∈ t ∧ ∀ j, z j ∈ (Ls j).lattice}

structure SubgroupData (Ls : Fin m → PeriodPair) where
  h0 : Submodule ℂ (Fin d0 → ℂ)
  U : Submodule ℚ (Fin d1 → ℚ)
  t : Submodule ℂ (Fin m → ℂ)
  h0_rat : RationalOverQbar h0
  t_ab : IsAbelianTangent Ls t
\end{Verbatim}

Here \lean{Lam} is $\Lambda(G)$ (the lattice disjunct is redundant, because Mathlib's $\wp$ takes the
junk value $0$ at lattice points, but it is clearer), \lean{Per} is $\Omega(G)$, and a
\lean{SubgroupData} describes the connected algebraic subgroup $G'=H_0\times T_U\times B$ with
tangent space $T_{G'}=h_0\times U^\perp\times t$: $h_0$ is spanned by algebraic vectors, $T_U$ is the
subtorus with Lie algebra $U^\perp$ (\texttt{torusLie\ U}), and $t$ is the tangent space of an abelian
subvariety $B=t/(t\cap\Lambda)$ of $E_1\times\dots\times E_m$. The codimensions are
$\delta_0=d_0-\dim h_0$, $\delta_1=\dim_\QQ U$, $\delta_2=m-\dim t$, $\delta=\delta_0+\delta_1+\delta_2$, and
\lean{Proper} means $\delta\ne0$. Ranks of subgroups are $\rk Y=\dim_\QQ\operatorname{span}_\QQ Y$. The
theorem is:

\begin{Verbatim}[fontsize=\scriptsize]
def Theorem41Abelian : Prop :=
  ∀ (d0 d1 m : ℕ) (Ls : Fin m → PeriodPair),
    (∀ j, IsAlgebraic ℚ (Ls j).g₂ ∧ IsAlgebraic ℚ (Ls j).g₃) →
    ∀ (V W : Submodule ℂ (Tan d0 d1 m)) (Y : AddSubgroup (Tan d0 d1 m)),
      finrank ℂ V < d0 + d1 + m → W ≤ V → RationalOverQbarTan W →
      Y.FG → (∀ y ∈ Y, y ∈ V ∧ y ∈ Lam Ls) →
      -- (NR)
      d1 + 2 * d0 + rk (Y ⊓ Per Ls) ≤ 2 * finrank ℂ V →
      ∃ H : SubgroupData (d0 := d0) (d1 := d1) Ls, H.Proper ∧
        finrank ℂ W - finrank ℂ ↥(W ⊓ H.tan) < H.δ ∧
        (((rk Y - rk (Y ⊓ (H.tan.toAddSubgroup ⊔ Per Ls))
              + H.δ1 + 2 * H.δ2 : ℕ) : ℤ) *
            ((d0 + d1 + m : ℕ) - finrank ℂ V : ℤ) ≤
          ((H.δ : ℤ) - (finrank ℂ W - finrank ℂ ↥(W ⊓ H.tan) : ℕ)) *
            ((d1 + 2 * m : ℕ) - rk (Y ⊓ Per Ls) : ℤ))

theorem LSG.W88.theorem41Abelian : Theorem41Abelian
\end{Verbatim}

Since $G'$ is connected, $\exp_G^{-1}(G')=T_{G'}+\Omega(G)$, so the first factor is $\lambda+\delta_1+2\delta_2$;
the second factor of the right side is $\delta-\tau$.

\begin{remark}[The subgroup family]
The theorem quantifies over products $H_0\times T_U\times B$. For
$\Ga^{d_0}\times\Gm^{d_1}\times A$ with $A$ an abelian variety, every connected algebraic subgroup is
such a product (Goursat, and the vanishing in characteristic $0$ of all Hom groups between additive,
multiplicative and abelian pieces), and a complex subtorus of $A$ is an abelian subvariety. So the
family is exactly Waldschmidt's. Inside the proof, the obstruction subgroup produced by the
multiplicity estimate is shown to have this form in Lean (Section~\ref{sec:subgroups}).
\end{remark}

\begin{remark}[Hypothesis \NR]\label{rem:NR}
\NR{} is not in \cite{W88}. It excludes, for instance, the standard encoding of the
Hermite--Lindemann theorem ($G=\Ga\times\Gm$, $V=W=\CC(1,1)$, $N=1$, where $1+2>2$); adding $\Gm$
factors restores \NR{} but makes the conclusion trivially satisfiable by a subtorus. All applications
in \cite{MstarPaper} satisfy \NR{} with room to spare.
\end{remark}

\subsection{Fidelity review}
An AI session that had not written the proof reviewed the statement against the page images of
\cite[pp.~385--386]{W88} and the formulation in \cite{MstarPaper}. Among other things it checked how
Lean elaborates each cast and natural-number subtraction (all $\NN$-subtractions are exact, and the
three integer subtractions $d-N$, $\delta-\tau$, $d_1+2m-\kappa$ are the intended ones) and that the
conclusion cannot be met trivially ($G'=G$ is excluded by \lean{Proper}, and $G'=0$ satisfies
\eqref{eq:W88} only when the expected inequality fails). Its verdict was \emph{correct}, with four
minor remarks: a docstring reference to a file outside the repository (since fixed), the need to
mention \NR{} in every citation, the redundancy of \lean{Proper} given $\tau<\delta$, and the
redundancy of the lattice disjunct in \lean{Lam}.

\subsection{Sanity derivations}
As evidence that the statement has teeth, \lean{LSG/W88/Sanity.lean} derives two classical theorems
from the \emph{proposition} \lean{Theorem41Abelian} and then instantiates them:
\lean{LSG.W88.Sanity.gelfondSchneider} (with $G=\Gm^2$, $V=W=\CC(1,b)$, $Y=\ZZ(x,bx)$, $\kappa=0$, where \NR{}
reads $2\le2$) and \lean{LSG.W88.Sanity.period_transcendental'}, Schneider's period theorem (with
$G=\Ga\times E$, $V=W=\CC(1,1)$, $Y=\ZZ(\omega,\omega)$). In both, a case analysis shows that every
admissible $H$ violates the conclusion. The second derivation reproves a result of Stage~1 by an
entirely different route.

% ===========================================================================
\section{Architecture}\label{sec:arch}
% ===========================================================================

Waldschmidt's proof \cite[\S\S6--7]{W88} has two inputs: an auxiliary polynomial $P$ that vanishes to
high order along $W$ at many points of $\exp_G(Y)$ with the saving $\kappa$ (Proposition~6.1, which
rests on \cite[Prop.~2.1, 2.4]{W83}), and Philippon's multiplicity estimate \cite{Philippon1986},
which turns such a $P$ into an obstruction subgroup. Section~7 of \cite{W88} combines them. The
formalization has three stages and a layer of bridges:
\begin{center}\small
\begin{tabular}{ccc}
\multicolumn{3}{c}{Stage 1: $\sigma,\zeta,\eta$, Legendre, addition theorem;
Schneider--Lang, Schneider (\S\ref{sec:stage1})}\\[2pt]
$\swarrow$ & & $\searrow$\\[2pt]
Stage 3a: W83 Prop.~2.1, Lemma U, & & Stage 2: multiplicity\\
W88 Prop.~6.1 (\S\ref{sec:aux}) & & estimate (\S\ref{sec:stage2}; StrongSix)\\[2pt]
$\searrow$ & & $\swarrow$\\[2pt]
\multicolumn{3}{c}{Stage 3b: W88 \S7 $\Rightarrow$ Theorem~\ref{thm:W88} under \NR{} (\S\ref{sec:sec7})}\\[2pt]
\multicolumn{3}{c}{$\downarrow$}\\[2pt]
\multicolumn{3}{c}{bridges: \S5d, abelian subvarieties, \lean{lean_mstar} (\S\ref{sec:5d})}
\end{tabular}
\end{center}
Stage~1 also feeds the multiplicity estimate (the group law of $E$ and its exponential) and Stage~3a
(quasi-periodicity of $\sigma$). StrongSix \cite{StrongSixRepo} is a Lake dependency: Stage~2 reuses
its filtered Hilbert functions, Bézout inequality, unmixedness, generic Wüstholz lemma and the
classification of algebraic subgroups of linear groups, and the bridges reuse its Theorem~1. Two
new block-generic files were added to StrongSix for this purpose; no existing StrongSix file was
changed.

Each piece was first stated as a \lean{Prop} interface (\lean{SchneiderLang},
\lean{WeierstrassOrderTwo}, \lean{MultEstAb}, \lean{Theorem41Abelian}) and the deductions between
interfaces were proved before their inputs, as in \cite{StrongSixPaper}. For example
\texttt{theorem41Abelian\_of\ :\ MultEstAb\ }$\to$\texttt{\ Theorem41Abelian} was complete before the multiplicity
estimate was; the main theorem is then one line, \texttt{theorem41Abelian\_of\ LSG.ME.multEstAb}.

% ===========================================================================
\section{Stage 1: elliptic functions and Schneider's theorems}\label{sec:stage1}
% ===========================================================================

\subsection{The Weierstrass library}
The directory \lean{LSG/Weierstrass/} (15 files, about 3{,}800 lines) defines, for a
\lean{PeriodPair}, the series
$\zeta(z)=1/z+\sum_{\omega\ne0}\bigl(1/(z-\omega)+1/\omega+z/\omega^2\bigr)$ and the product
$\sigma(z)=z\prod_{\omega\ne0}(1-z/\omega)\exp(z/\omega+z^2/2\omega^2)$. In Lean the factor for $\omega=0$
evaluates to $1$, so the product runs over the whole lattice. The construction uses Mathlib's API for
infinite products of holomorphic functions (developed there for the Dedekind $\eta$ function). The main results are: $\sigma$ is entire, vanishes exactly on
$\Lambda$, is odd with $\sigma'(0)=1$, and $\sigma'/\sigma=\zeta$; $\zeta'=-\wp$; the quasi-periods
$\eta(\omega)=\zeta(z+\omega)-\zeta(z)$ are $\ZZ$-linear on $\Lambda$; $\sigma(z+\omega)=\pm
e^{\eta(\omega)(z+\omega/2)}\sigma(z)$ (\lean{sigma_add}); $\sigma$ has order two, and $\wp,\wp'$ are
quotients of entire functions of order two by $\sigma^3$ (\lean{weierstrassOrderTwo}).

\emph{Legendre's relation} (\lean{LSG.Weierstrass.legendre}) states
$\eta(\omega_1)\omega_2-\eta(\omega_2)\omega_1=\pm2\pi i$, with the sign of $\operatorname{Im}(\omega_2/\omega_1)$
(Mathlib's \lean{PeriodPair} carries no orientation). Mathlib has Cauchy's theorem for rectangles but
no residue theorem for parallelograms. We use only rectangles: for a basis $(\alpha,\beta)$ with
$\tau=\beta/\alpha$ in the upper half-plane, $f(z)=\alpha\zeta(\alpha z)-\alpha\eta(\alpha)z$ is $1$-periodic and
shifts by a constant under $z\mapsto z+\tau$, and $f=G+1/z$ with $G$ holomorphic near the rectangle
$[-\frac12,\frac12]\times[-\frac{\operatorname{Im}\tau}2,\frac{\operatorname{Im}\tau}2]$. The vertical sides
cancel, the horizontal sides differ by the constant, $\oint G=0$ by Cauchy--Goursat, and $\oint
dz/z=2\pi i$ is computed with explicit branches of $\log$ (\lean{rectIntegral_inv}). The same trick
gives the ``sum of residues'' fact we need (an elliptic function with at most one simple pole per
period is constant) by subtracting $r\zeta(z-a)$ and using Legendre.

\emph{The addition theorem} (\lean{weierstrassP_add}) and the \emph{degree-two property}
($\wp(z)=\wp(w)$ iff $z\equiv\pm w$, \lean{weierstrassP_eq_iff}) give the \emph{dictionary}
(\lean{LSG/Weierstrass/Dictionary.lean}): $\exp_E(z)=(\wp(z),\wp'(z)/2)$ is a homomorphism from $\CC$ to
the group of points of Mathlib's Weierstrass curve $y^2=x^3-\frac{g_2}4x-\frac{g_3}4$ with its
algebraic group law (\lean{expE_add}). Consequences used later: the discriminant is nonzero, and
points with coordinates in a field $K\ni g_2,g_3$ form a subgroup.

\subsection{The Schneider--Lang criterion and Schneider's theorems}
The proposition \lean{LSG.SchneiderLang} (\lean{LSG/Interface.lean}) is the one-variable criterion of
\cite[Ch.~III, Thm.~1]{Lang1966} for meromorphic functions: if $f_i=g_i/h$ with $g_i,h$ entire of order
$\le\rho$, each $f_i'$ is a polynomial over a number field $K$ in the $f_j$ on $\{h\ne0\}$, and two of the
$f_i$ are algebraically independent over $K$, then only finitely many $w$ with $h(w)\ne0$ have all
$f_i(w)\in K$. The proof \lean{LSG.schneiderLang} (\lean{LSG/SchneiderLang/}, about 1{,}500 lines)
is the classical one: Siegel's lemma over $\mathcal O_K$ (Mathlib's
\lean{house.exists_ne_zero_int_vec_house_le}), a Schwarz lemma with multiplicities, Cauchy
estimates, and Liouville's inequality; it shows that the number of such points is at most
$4[K:\QQ]\lceil\rho\rceil$.

Schneider's theorems follow by taking $h=\sigma^3$: with $e=\wp(\omega/2)$ for a primitive period
$\omega$, the functions $z,\wp,\wp'$ take values in $K=\QQ(g_2,g_3,e,\omega)$ at the infinitely many points
$(k+\frac12)\omega$ if $\omega$ were algebraic (\lean{LSG.Schneider.period_transcendental}); the functions
$\wp,\wp',e^{2\pi iz/\omega}$ do the same if $\omega/\pi$ were algebraic
(\lean{period_div_pi_transcendental}); and translates of an algebraic point give
\lean{wp_transcendental}.

\subsection{The lemniscatic invariant}
For \lean{lean_mstar} we need $g_2(\vp\ZZ[i])=4$ with $\vp=\Gamma(1/4)^2/(2\sqrt{2\pi})$
(\texttt{LSG.LemniscateIndep.g}$_2$\texttt{\_eq\_four\_of}, about 600 lines). The proof uses the symmetry $z\mapsto iz$
($g_3=0$), the real branch of $\wp$ on $(0,\vp/2]$ with the substitution $t=\sqrt{e/\wp}$, and the
Beta integral for $2\int_0^1dt/\sqrt{1-t^4}$. With Schneider's theorem this gives
$\vp/\pi\notin\Qbar$ (\lean{LSG.Bridge.irrational_varpi_div_pi_of}). See Section~\ref{sec:trust} for
the history of this part.

% ===========================================================================
\section{Stage 2: the multiplicity estimate with elliptic factors}\label{sec:stage2}
% ===========================================================================

\subsection{Statement}
\lean{LSG.ME.MultEstAb} (\lean{LSG/ME/MultEstStmt.lean}) is Philippon's multiplicity estimate
\cite[Thm.~2.1]{Philippon1986} on $G=\Ga^{d_0}\times\Gm^{d_1}\times E_1\times\dots\times E_m$ over $\CC$, in
the shape of StrongSix's \lean{MultEst}. Polynomials live in $\CC[X;Y,Z;(U_k,V_k,W_k)_{k\le m}]$, with
$E_k\subset\PP^2$ in each elliptic block and $Y_jZ_j=1$ for the torus. Let $n=d_0+d_1+m$. If
$P\notin I_G$ has block degrees $\le(D_0,D_1)$, is homogeneous of degree $D_2$ in each elliptic block,
and has formal order $>nT$ along $W$ at every point of $\Sigma(n)=\{\sigma_1+\dots+\sigma_n:\sigma_i\in S\}$
($0\in S$), then there is a proper obstruction $G'$ such that every $F\subset S$ of points pairwise
inequivalent modulo $G'$ satisfies
\[
T^{\,\dim W-\dim(W\cap T_{G'})}\,|F|\,D_0^{\dim h_0}D_1^{\,d_1-\dim U}D_2^{\dim t}\le C\,D_0^{d_0}D_1^{d_1}D_2^{m}.
\]
The obstruction is an \lean{ObstrData}: tangent data $(h_0,U,t)$ with $t$ an abelian tangent space,
together with the subgroup $\overline{\exp_G(T_{G'})}^{\mathrm{Zar}}$ and the fact that its Lie algebra is
$T_{G'}$. The proof \lean{LSG.ME.multEstAb} produces $C=n^nd_0^{d_0}d_1^{d_1}4^m2^nn!\,4^{d_0}8^{d_1}12^m$.
$\Qbar$-rationality of $h_0$ is not part of the output; it is restored in \S\ref{sec:rat}.

StrongSix follows Roy's affine treatment \cite[Ch.~5, 8]{DALAG}, which works for linear groups because
their invariant derivations and translations preserve degrees. Neither holds with elliptic factors.
We kept Roy's architecture (filtered Hilbert functions, Bézout by complete-intersection chains,
Wüstholz's lemma by linear algebra, prime ideals instead of components) and changed four things.

\subsection{The model and the addition laws}
Each $E_k$ is the cone over $V^2W=4U^3-g_2UW^2-g_3W^3$ in $\CC^3$; the ideal $I_G=(Y_jZ_j-1,F_k)$ is a
prime complete intersection. Translations use a complete system of two addition laws of bidegree
$(2,2)$ on $E\subset\PP^2$ (Mathlib's projective addition formula and a Bosma--Lenstra law
\cite{BosmaLenstra1995}), proved in Lean to preserve the curve equation, to represent the sum whenever
nonzero, and never to vanish simultaneously (\lean{LSG/Elliptic/AdditionLaws*.lean}). Translation then
multiplies elliptic block degrees by $2$; this is Philippon's constant $c_1$.

\subsection{Degree growth of derivations: Philippon's operators}
Roy's main argument builds a chain $L_r=I_G+(\tau_\sigma D_{\mathbf l}P)$ from derivation words $D_{\mathbf
l}$ of length $\le rT$. On $E$ the invariant derivation, written on the cone in a chart, raises the
degree by one, so the generators of $L_r$ have elliptic degree about $D_2+rT$, and Bézout then bounds
everything in degree $D_2+nT$ instead of $D_2$. In Waldschmidt's parameters $T=\Delta/\log S$ and
$D_2=\Delta/S^2$, so $T\gg D_2$: the estimate obtained this way is useless. (A first version of the chain
was formalized and then discarded for this reason.)

The fix is Philippon's own device \cite[\S4.1, Déf.~4.1--4.2, Prop.~4.3--4.4]{Philippon1986}: instead of
iterating derivations, differentiate along a curve of fixed degree,
\[
\delta^\kappa_gQ=\partial_z^\kappa\,Q\bigl(\mathrm{law}(X,\mathrm{law}(\mathrm{rep}\,g,\,\exp(zw)))\bigr)\big|_{z=0}.
\]
The point $X+g+\exp(zw)$ is represented by polynomials in $X$ of fixed degree with coefficients that are
formal power series in $z$, so every $\delta^\kappa_gQ$ has elliptic degree $\le2D$ for all $\kappa$
(\lean{phOp_mem_FLE}). Algebraically, $\exp(zw)$ is the formal exponential: $x+zw$ on $\Ga$, $y\exp(zw)$
on $\Gm$, and on $E$ the formal solution $(s(z):1:t(z))$ of the ODE in the $V$-chart,
$s'=-\frac12-g_2st-\frac32g_3t^2$, $t'=-6s^2+\frac{g_2}2t^2$, $s(0)=t(0)=0$, with coefficients in
$\QQ[g_2,g_3]$. Proving Philippon's Proposition~4.4 (zeros of $\delta^\kappa_gQ$) needs the formal
associativity of the laws along such curves (a uniqueness theorem for formal ODE solutions) and an
identification of the formal order with vanishing of derivation words in canonical charts. Derivation
words survive only inside Wüstholz's lemma, where they shift the Hilbert function argument by a
constant, which does not affect leading terms.

\subsection{Cone vertices: affine slices}\label{sec:slices}
StrongSix works with affine (filtered) Hilbert functions on $K[X,Y,Z]$, where there are no vertices.
On the cone model every generator of $L_r$ and of Roy's ideal $A$ is homogeneous of positive degree in
each elliptic block. Hence the locus $\{(U_k,V_k,W_k)=0\}$, of dimension $d_0+d_1+2m-2$, lies in the
affine zero set of every $L_r$, the pigeonhole on Hilbert dimensions becomes vacuous, a prime of
maximal dimension may have no projective zeros at all, and the affine Nullstellensatz fails. In
Philippon's multiprojective setting these irrelevant primes have zero Hilbert function; in the affine
setting they dominate.

We separate the two uses of ideals. For geometry we pass to \emph{relevant radicals}
$I(Z(L))$ (only points with $(U_k,V_k,W_k)\ne0$ in every block): the pigeonhole, the choice of a
relevant prime and Roy's component step are done there (\lean{MainAbChain.lean}). For Hilbert functions
we \emph{slice}: choose $\alpha\in\CC^m$ and replace an ideal $I$ by $I+(\ell_k-1)_k$ with
$\ell_k=V_k+\alpha_kW_k$. For generic $\alpha$ (avoiding finitely many points, \lean{exists_good_alpha}) the
slices of the relevant primes keep their dimension minus $m$, the sliced $I_G$ is a complete
intersection with an explicit standard-monomial bound for its Hilbert function (\lean{SliceIG.lean}),
and there is a slice Nullstellensatz $\sqrt{I+(\ell_k-1)}=I(Z_\alpha(I))$ (\lean{radical_sl}). Primality of
partial slices is proved by homogenizing with powers of $\ell_k$, without localization. Bézout and the
final count are then carried out for sliced ideals.

\subsection{Projective Wüstholz: a multiplier and glued charts}
Roy's form of Wüstholz's lemma \cite[Lemma~8.16, Thm.~8.18]{DALAG} needs derivations $D_i$ and
polynomials $Q_j$ in the ideal of the coset union with $D_iQ_j\equiv\delta_{ij}$. On a projective
model one only gets $D_iQ_j\equiv\delta_{ij}M$ with $M$ a non-zero-divisor (\lean{wustholzM'}). Moreover a
single chart does not work for all cosets: on cosets whose points have a constant $2$-torsion
coordinate ($V=0$) the $V$-chart derivation vanishes. We glue chart-dependent derivations with
separating polynomials (\lean{glue}). The sliced version needs a multihomogeneous $M$ (so that
\lean{nzd_sl}: $M$ is a non-zero-divisor modulo the sliced ideal), and Euler-corrected derivations
$D_i'=D_i-\sum_k(D_i\ell_k)\,\mathrm{Eul}_k$ that preserve the slice (\lean{wustholz_cosets_sl}).

\subsection{The obstruction subgroup}\label{sec:subgroups}
The stabilizer $H=\{g: g+Z(p)\subseteq Z(p)\}$ of the zero set of a relevant prime is a Zariski-closed
subgroup (a Noetherian argument for any ideal, \lean{StabG.lean}). To turn it into data
$(h_0,U,t)$ we proved:
\begin{itemize}
\item $\exp_G^{-1}(H)$ is a closed subgroup of $T_G$ containing $\Omega(G)$; its real linear part is a
complex subspace \lean{lieAlg} (identity theorem), and its image in $\CC^m$ meets $\Lambda$ in a lattice.
\item The projection of $H$ to $\Ga^{d_0}\times\Gm^{d_1}$ is Zariski closed (projective elimination along
$(\PP^2)^m$, \lean{isClosed_proj}), hence of the form classified in StrongSix.
\item \emph{Zariski tangent space $=$ Lie algebra} (\lean{mem_TZ_iff}). The naive argument (``$I(H)$ is
translation invariant, hence $D_v$-stable'') is false as written on the cone, since $D_vP$ is not
multihomogeneous (a counterexample: the diagonal of $E\times E$ with $P=U_1W_2-U_2W_1$); the proof
homogenizes $D_v$ with chart forms.
\item \emph{Dimension.} The algebraic dimension of $H^0=\overline{\exp_G(\lieAlg)}$ equals
$\dim\lieAlg+m$ (\lean{hdim_H0}), via a bridge between StrongSix's Hilbert dimension and transcendence
degree (\lean{hdim_eq_trdeg}) and the existence of smooth points in characteristic $0$
(\lean{exists_smooth_point}): $\le$ at the identity from the tangent space, $\ge$ at a smooth point.
\item \emph{Product decomposition} (\lean{lieAlg_eq_prod}): $\lieAlg=h_0\times U^\perp\times t$. First,
the abelian projection of $\lieAlg$ is $\Lambda$-rational, by a transcendence-degree count through the
exponential parametrization. Second, a graph part would give a $\CC$-linear form on that rational space
with purely imaginary values ($2\pi i\ZZ$ or $0$) on its lattice points, hence zero.
\end{itemize}
Purely analytic arguments cannot replace the algebraic inputs here: the graph of the Tate
uniformization $\CC^\times\to E(\CC)$ is a closed complex Lie subgroup of $\Gm\times E$ that is not algebraic,
so no argument using only the closed-subgroup structure of $\exp_G^{-1}(H)$ can yield the product
decomposition. The second step above uses that $H^0\cap(\Ga^{d_0}\times\Gm^{d_1})$ has finitely many
components, a consequence of Noetherianity.

The Hilbert function lower bound for $H^0$ (degree $t_0+t_1+t_2+m$, the extra $m$ from the free cone
scalings in blocks not met by $t$) and the coset count finish the subgroup layer. The assembly
(\lean{LSG/ME/MainAb.lean}) follows Roy \S8.4.1 with the slicing of \S\ref{sec:slices}: chain and relevant
pigeonhole, relevant prime $p$, stabilizer and $H^0$, Roy's ideal $A$, component step, a good $\alpha$ for all
cosets, Wüstholz for the sliced ideals, Bézout for the sliced $A$, and the final arithmetic.

% ===========================================================================
\section{Stage 3a: the auxiliary function with periods}\label{sec:aux}
% ===========================================================================

\subsection{W83 Proposition 2.1 in our setting}
The $\kappa$-saving comes from \cite[Prop.~2.1]{W83}: Siegel's lemma is applied with conditions only
on a fundamental domain of the periods $\omega_1,\dots,\omega_\kappa$ in a ball of radius $r$, which saves
a factor $r^{\kappa}$ in the number of conditions; quasi-periodicity of the auxiliary function then
propagates smallness to the whole ball, and an interpolation lemma converts smallness on a point set
into smallness on the ball. We formalized this with two simplifications.
\begin{enumerate}
\item In all our applications the periods of $G$ lie over periods of the theta function, so the
multiplier $\Psi(z+\omega)/\Psi(z)$ is an explicit exponential given by \lean{sigma_add}. The
lower bound for the theta function near the points (\cite[Lemma~2.2]{W83}) is therefore not needed
inside Prop.~2.1.
\item Waldschmidt's interpolation \cite[Lemma~2.3]{W83} rests on Masser's interpolation theorem for
separated sets in $\CC^n$. We avoid it by choosing coordinates of $V$ among the coordinate forms of
$T_G$: these send $\Omega(G)$ into $0$, $2\pi i\ZZ$ or a lattice $\Lambda_j$, so the grid can be taken to be
a product of one- and two-dimensional progressions, and interpolation is iterated one-variable Lagrange
interpolation combined with StrongSix's one-variable Schwarz lemma. For two-dimensional grids the node
products are bounded through $|u\alpha+v\beta|\ge c\max(|u|,|v|)$ and $\max(|u|,|v|)^2\ge\max(|u|,1)\max(|v|,1)$,
and the number of period classes in the grid is counted with an adjugate matrix instead of a Smith
normal form.
\end{enumerate}
The result \lean{prop21_G} is W83's Prop.~2.1 on a subspace $V\subset T_G$, with one change: the
hypothesis $U\ge r^2\log(R/r)$ becomes $c_1r^2\log(R/r)\le U$, because the coordinates $\CC^n\cong V$ rescale
radii. In the application ($r=\nu^2S$, $R=\nu^3S$, $U=\nu^6\Delta$) the constant is absorbed by $\nu$.

\subsection{Lemma U and Proposition 6.1}
Proposition~6.1 of \cite{W88} asks for vanishing \emph{along $W$} at the points of $\exp_G(Y(S))$, $Y(S)=\{\sum
h_jy_j:0\le h_j\le S\}$. Following Lemma~U of \cite{MstarPaper}, the small values of $Q_S$ on a ball give,
through Cauchy's estimates, small Taylor coefficients along $W$ at each point, and Liouville's
inequality makes them vanish. Our formal proof (\lean{lemmaU_core}) differs from that sketch in ways
that simplify it.
\begin{itemize}
\item \emph{A chart per point.} Instead of a fixed theta coordinate and a decomposition $Y=\mathcal T+F$
with translations by $F$, for every $u\in\CC$ some coordinate of $(\sigma^3,\sigma^3\wp,\sigma^3\wp')(u+z)$ is
$\ge e^{-c(1+|u|)^2}$ for $|z|\le\rho_0$ (\lean{exists_chart}, by quasi-periodicity and compactness of a
fundamental domain). So the argument runs at every point of $Y(S)$ and the polynomial is $Q_S$ itself.
\item \emph{A fixed base point.} The second argument of the addition law is $(\wp,\wp',1)(z_0+u)$ with
$z_0=\omega_1/2$, whose Taylor coefficients satisfy the recursion of $\wp''=6\wp^2-g_2/2$ with algebraic
initial values; no derivatives are taken at the point $\exp_G(t)$ itself.
\item \emph{One variable.} Vanishing along $W$ is reduced to finitely many lines $s\mapsto t+sw$ and
recovered by polarization.
\item \emph{Sizes} are house and denominator in an abstract number field. Algebraic points
$\exp_G(\sum h_jy_j)$ are computed by balanced splitting with whichever addition law is nonzero, which
keeps their sizes $\le e^{cS^2}$.
\end{itemize}
\NR{} enters exactly once, to guarantee $\Delta\ge S^2$, which gives the hypothesis $c_1r^2\log(R/r)\le U$ of
Prop.~2.1. The result \lean{LSG.W88.prop61_b} is Prop.~6.1 with $a=1$ and any integer $b\ge1$
(Waldschmidt's Remark~B: $D_0(\log S)^b=\Delta$, $\Delta^{d-N}=CS^{d_1+2m-\kappa}(\log S)^{bd_0}$), with the auxiliary
function nonzero on $G$ and vanishing to order $>\Delta/\log S$ along $W$ on $Y(S)$.

% ===========================================================================
\section{Stage 3b: Waldschmidt's \S7}\label{sec:sec7}
% ===========================================================================

\texttt{theorem41Abelian\_of\ :\ MultEstAb\ }$\to$\texttt{\ Theorem41Abelian} (\lean{LSG/W88/Sec7*.lean}, about 1{,}700
lines) follows \cite[pp.~396--397]{W88}.

\emph{From analytic to formal vanishing.} Prop.~6.1 gives analytic vanishing of $\Psi=\sum p_l M_l\circ\exp_G$
along $W$; the multiplicity estimate wants formal order of the polynomial $P$ in Philippon's sense.
The bridge \lean{ordGt_of_vanishAlong} works only with Taylor coefficients at one point: the Taylor series
of the analytic lift of $s\mapsto\exp_G(y+sw)$ in the canonical chart is a formal flow with the right
initial value, hence equals the formal exponential by uniqueness of formal ODE solutions
(\lean{tps_eq_fflow}); a unit factor relates $\Psi$ and $P$ along the line.

\emph{Counting.} If $u\in\exp_G^{-1}(G')$ then $Nu\in T_{G'}+\Omega(G)$ for some $N\ge1$ (\lean{htil_le_span}),
using only the fields of \lean{ObstrData}. Hence $(M+1)^\lambda$ points of the grid are pairwise
inequivalent modulo $G'$ (\lean{exists_inequiv_grid}).

\emph{The inequality.} With Prop.~6.1 at $b=2$ and the multiplicity estimate at
$T'=\lfloor\Delta/(2n\log S)\rfloor$, $D_0=\lfloor\Delta/\log^2S\rfloor+1$, $D_1=\lfloor\Delta/S\rfloor+1$, $D_2=D+1$
(the polynomial is multiplied by $\prod_kW_k$ so that $D_2\ge1$), taking logarithms gives $A\log S\le
K+B\log\log S$ with $A=q(\lambda+\delta_1+2\delta_2)-e(\delta-\tau)$, $B=2d_0(\delta-\tau)+q(\tau-2\delta_0)$, $q=d-N$,
$e=d_1+2m-\kappa$, and $A,B,K$ bounded independently of $S$. One large $S$ gives $A\le0$, which is
\eqref{eq:W88}; if $\delta\le\tau$ one finds $\delta=\tau=\delta_0\ge1$ and $B=-q\delta_0<0$, a contradiction. No
pigeonhole over $S$ is needed.

\subsection{The $\Qbar$-rationality trick}\label{sec:rat}
The obstruction of \lean{MultEstAb} lives over $\CC$, and its additive part $h_0$ need not be
$\Qbar$-rational: a subgroup like $\CC\cdot(1,\pi)\subset\Ga^2$ is a perfectly good closed subgroup over $\CC$.
W88 requires $G'$ defined over $\Qbar$. Rather than descend the stabilizer to $\Qbar$, we replace $h_0$ by
its algebraic part $h_0''=\operatorname{span}_\CC\{a\in h_0: a\text{ algebraic}\}$ (\lean{exists_subgroupData_rat}).
This does not change $\lambda$, because the $\Ga$-coordinates of points of $Y$ are algebraic, so $y\in
T_{G'}+\Omega$ implies $y\in T_{G''}+\Omega$; it leaves $\delta_1,\delta_2$ unchanged; and $\delta-\tau$ does not decrease,
since $\dim(W\cap T'')\ge\dim(W\cap T)-(\dim h_0-\dim h_0'')$. As $d_1+2m-\kappa>0$ under \NR, both $\delta>\tau$ and
\eqref{eq:W88} transfer.

% ===========================================================================
\section{The \S5d bridges}\label{sec:5d}
% ===========================================================================

Applications often use Waldschmidt's hyperplane consequence \cite[\S5d]{W88}: if $V$ is a hyperplane,
$Y$ has rank $m$ and $m>(d_1+2d_2-\kappa)(d-1-t)$ ((5.4); $t=\dim W$), then $V$ contains $T_{G'}$ for some
connected algebraic subgroup $G'\ne0$. Waldschmidt's proof chooses a subgroup $G'$ maximizing a ratio
$a(G')/e(G')$ (Lemma~5.1), applies Theorem~\ref{thm:W88} to $G'$ with $V'=V\cap T_{G'}$, and concludes with
Lemma~5.2. Two difficulties appear in the formalization.

\emph{Theorem~\ref{thm:W88} on the maximizer.} $G'$ is $\Ga^{k_0}\times\Gm^{k_1}\times B$ with $B$ an abelian
subvariety, not a product of curves, and \NR{} may fail for $G'$ even if it holds for $G$. We proved
\lean{LSG.CM.thm41Sub}: Theorem~\ref{thm:W88} under \NR{} for $\Ga^{k_0}\times\Gm^{k_1}\times B$, $B$ any
abelian subvariety of $E_1\times\dots\times E_m$. No complex multiplication is needed: for a suitable
choice of $\dim B$ coordinates, the coordinate projection $B\to\prod_kE_{\iota(k)}$ is an isogeny, and the
theorem transports along its tangent map (logarithms pull back to logarithms, periods have torsion
quotient, subgroups push forward with the same $\delta$, $\tau$). When \NR{} fails for $G'$ there are two
cases. If $B=0$, $G'$ is linear and StrongSix's Theorem~1 (Waldschmidt's theorem for linear groups, in
Roy's form, without \NR) applies after converting Roy's surjection into a subgroup
(\lean{sub41Concl_bot}). If $B\ne0$, failure of \NR{} means $\kappa'\ge d_1'+2d_2'-1$, and the periods in $V$
then force $T_B\subset V$, which contradicts the assumption.

\emph{Rationality of abelian tangent spaces.} The maximizer must be defined over $\Qbar$; for products of
copies of one curve (with or without CM) we prove that abelian tangent spaces are $\Qbar$-rational. For
different curves this is assumed as a hypothesis \lean{hrat} (\lean{LSG.W88.hypStd_of_rat}); it is true
for curves over $\Qbar$, but not proved here.

The resulting endpoints are \lean{LSG.W88.thm11Kappa} ($\Gm^{d_1}\times E$),
\lean{LSG.W88.hypStdCM} and \lean{hypStdCMAt} ($\Ga^{d_0}\times\Gm^{d_1}\times E^r$). Their statements are
byte-for-byte copies of \lean{lean_mstar}'s hypotheses (\lean{LSG/W88/MStarCopy*.lean}), and we checked
that, e.g., \texttt{exact\ LSG.W88.thm11Kappa\_of\_theorem41\ h\ L\ d}$_1$ typechecks against \lean{lean_mstar}'s
unmodified definitions.

% ===========================================================================
\section{Vector extensions}\label{sec:vecext}
% ===========================================================================

\subsection{The class and the statements}
The applications MG$^\natural$, MX, MX$'$ and the four exponentials statement at $t=c\vp$ of
\cite{MstarPaper} use groups built from universal vector extensions $E^\natural$ (extensions of $E$ by
$\Ga$, with periods $(\lambda,\eta(\lambda))$ for $\lambda\in\Lambda$). We work with the class of groups
\[
G=\Ga^{d_0}\times\Gm^{d_1}\times G_2,\qquad G_2=\bigl((E_1\times\dots\times E_m)^\natural\times\Ga^r\bigr)\big/\{(-u,Mu)\},
\]
the push-out of $(E_1\times\dots\times E_m)^\natural$ along an algebraic matrix $M\colon\CC^m\to\CC^r$, so
$d_2=\dim G_2=m+r$. On the tangent space $T_G=\CC^{d_0}\times\CC^{d_1}\times\CC^m\times\CC^r$ the
exponential of $E^\natural$ is $\exp(z,w)=((\wp_j(z_j),\wp_j'(z_j)),\,w-\zeta(z))$, with $s=Mw$ on the last
factor. Hence $\Lambda(G)$ is described by ``$s-M\zeta(z)$ algebraic'' (with Lean's $\zeta$, which takes the
value $\eta(\lambda)$ at $\lambda\in\Lambda$), and $\Omega(G)=0\times2\pi i\ZZ^{d_1}\times\{(\lambda,M\eta(\lambda))\}$. The
\emph{fibre} is the vector part $0\times0\times0\times\CC^r$ of $G_2$. Connected algebraic subgroups are
$T_U\times(\tilde B+\Gamma)$, with $\tilde B$ the image of $B^\natural$ for an abelian subvariety $B$ and $\Gamma$ a
$\Qbar$-rational subspace of $\Ga^{d_0}\times\text{fibre}$; they are not products. The statements are in
\lean{LSG/W88/StatementVec.lean}, \lean{V5d2Stmt.lean} and \lean{V5dOnto.lean}; their conclusion is that of
\lean{Theorem41Abelian} with $d_2=m+r$. The general statement \lean{Theorem41Vec} has no further premises.

Two premises concern the curves. \lean{LiftAlg}: if $z\in T_B$ and $\exp(z)\in B(\Qbar)$, there is $w$ with
$(z,w)\in T(B^\natural)$ and $w-\zeta(z)$ algebraic. \lean{PerTanRat}: $T(B^\natural)$ is $\Qbar$-rational.
Both are true for all curves with algebraic invariants ($B^\natural\to B$ is a surjective morphism defined
over $\Qbar$), but we have proved them in Lean only for the CM curves used by the applications
(\lean{LSG.CM.liftAlg_block}, \lean{perTanRat_block}, and the products $E_i^n\times E_\rho^n$ of the MX
bridges).

\begin{theorem}[Lean; standard axioms]\label{thm:vec}
Let $G$ be in the class above, with $g_2(\Lambda_j),g_3(\Lambda_j)$ and $M$ algebraic.
\begin{enumerate}
\item[(a)] If $r\le1$, $V\cap\text{fibre}=0$ and the curves satisfy \lean{LiftAlg}, Theorem~\ref{thm:W88}
holds under \NR{} (\lean{LSG.W88.theorem41VecLA}, for every algebraic $M$; the case of onto $M$,
\lean{theorem41VecLA_onto}, was proved first, by a shorter route).
\item[(b)] If $d_0=0$, $r\le2$, $\dim(V\cap\text{fibre})\le1$ and the curves satisfy \lean{LiftAlg} and
\lean{PerTanRat}, Theorem~\ref{thm:W88} holds under \NR{} (\lean{LSG.W88.theorem41Vec2}).
\item[(c)] The \S5d consequence ((5.4) with $d_2=m+r$, without \NR) holds for $r\le1$, and for $d_0=0$,
$r\le2$, for curves satisfying \lean{LiftAlg} and \lean{PerTanRat}
(\lean{LSG.W88.V5d.hyperplaneKappaVec_uncond}, \lean{hyperplaneKappaVec2_uncond}).
\end{enumerate}
\end{theorem}

The restrictions on $r$, $d_0$ and $V\cap\text{fibre}$ are genuine; the general \lean{Theorem41Vec} is not
proved. Accordingly these results are to be cited as ``W88 Thm.~4.1 under \NR{} for
$\Ga^{d_0}\times\Gm^{d_1}\times G_2$, $G_2$ a push-out of $(\prod E_j)^\natural$ along an algebraic $M$, in the
cases $r\le1$ with $V\cap\text{fibre}=0$, or $d_0=0$, $r\le2$ and $\dim(V\cap\text{fibre})\le1$; for curves
satisfying \lean{LiftAlg} (and \lean{PerTanRat}), true for all curves over $\Qbar$ and Lean-proved for the
CM curves used'', and never as ``W88 Thm.~4.1 for vector extensions''.

\subsection{The proof}
The architecture is that of Sections~\ref{sec:stage2}--\ref{sec:sec7}.
\begin{itemize}
\item \emph{A model of $E^\natural$.} The surface $X=V(F,\,WT_3-VT_1+2U^2)\subset\PP^4$, with $(U:V:W:T_1:T_3)=(x:y:1:t:yt-2x^2)$
over affine points, is $E^\natural$ plus a boundary line. Exact linear algebra over $\mathbb F_p$ shows that it
carries no addition law of bidegree $(2,2)$ (the missing section is $Ut$); we found eight laws of bidegree
$(3,3)$ and proved in Lean that they preserve the ideal, represent the sum in all configurations (and
analytically, \lean{NLaw_expN_all}), and form a complete system (\lean{NLaw_complete}). The ideal $(F,G)$ is
prime (\lean{isPrime_span_FG}).
\item \emph{The multiplicity estimate} \lean{LSG.MEV.multEstVec} (\lean{LSG/MEV/}). Every $G$ of the class is
the quotient of $\tilde G_r=\Ga^{d_0+r}\times\Gm^{d_1}\times E_1^\natural\times\dots\times E_m^\natural$ by the vector
subgroup $K_M=\{(0,-Mv;0;0,v)\}$, and polynomials on $G$ are $K_M$-invariant polynomials on $\tilde G_r$. All
layers of Section~\ref{sec:stage2} were redone for $E^\natural$ blocks: primality and the complete
intersection property of $I_{\tilde G}$, the Nullstellensatz despite the boundary plane, translations, the
formal exponential and Philippon's operators (degree $\le3D$), the chain of ideals with its relevant-prime
step, the $E^\natural$ curve lemma, the subgroup layer (unconditional classification of Zariski-closed
subgroups, $\hdim H^0$, finiteness of cosets, and a Hilbert lower bound matching $\hdim H^0$), the slice
layer, and Wüstholz's lemma (homogenized chart derivations, transversals at every coset, coset bounds). The
cone over $X$ is singular along $U=V=W=0$, yet the change of chart for the homogenized derivations is
polynomial, with explicit cofactors. When $M$ is onto, the grading of Section~\ref{sec:stage2} (all additive
coordinates in one block) gives the estimate. When it is not (for instance $r=1$, $M=0$, so that
$G_2=E^m\times\Ga$, a case that W88~\S7 needs), the trailing additive coordinates belong to $G_2$ and must
carry the elliptic degree $D_2$; we use a grading with four blocks (leading and trailing additive
coordinates, torus, elliptic). Hilbert \emph{dimensions} do not depend on the grading, so only the degree
steps had to be redone (about 2{,}200 lines).
\item \emph{The analytic side}: Prop.~6.1 for the class (\lean{LSG.AuxV.prop61V_b}), using entire
coordinates of $\exp_{E^\natural}$ built from $\sigma$ and the premise $V\cap\text{fibre}=0$ to keep the grid
coordinates discrete (the fibre forms have non-discrete period images when $m\ge3$). Sizes of fibre values
are controlled through the $\zeta$ addition formula rather than the $(3,3)$ laws, whose balanced splitting would
give sizes $6^k\approx S^{2.58}$ and break the $e^{cS^2}$ bounds. For case~(b), where $V\cap\text{fibre}=\CC u$
may be an irrational line, we proved W83's Prop.~2.1 with one polynomial direction (\lean{prop21LatT}): the
fibre forms are affine in $s$, so along $u$ the functions are polynomials of degree $\le D$; the Siegel
conditions are imposed at $D+1$ nodes in that variable and recovered by exact Lagrange interpolation, at the
cost of one power of the parameter $U$.
\item \emph{W88 \S7 and $\Qbar$-rationality.} The obstruction lives over $\CC$, with tangent space
$T'=\Gamma_{\mathrm{full}}\oplus U^\perp\oplus T(\tilde B)$. For $r\le1$ the fibre part of $\Gamma_{\mathrm{full}}$ is $0$ or the
whole fibre, hence rational, and replacing $\Gamma_{\mathrm{full}}$ by the span of its algebraic vectors works as in
Section~\ref{sec:rat}; \lean{LiftAlg} keeps $\lambda$ unchanged. For $d_0=0$, $r\le2$, $\Gamma_{\mathrm{full}}$ may be an
irrational line in the two-dimensional fibre. Then $T'$ lies between a rational subgroup $T_S$ (its algebraic
part) and a rational subgroup $T_E$ (the whole fibre added) with $\dim T_E/T'\le1$, and one of the two
satisfies the inequality \eqref{eq:W88}: if $T_E$ fails, \lean{PerTanRat} puts every algebraic vector of
$T'$ into $T_S$, and $T_S$ works.
\item \emph{The \S5d engine} follows Waldschmidt's Lemmas~5.1--5.2 as in Section~\ref{sec:5d}, with frames
for the class. For $r\le1$ the maximizing subgroup may have a non-onto push-out matrix; moving its $\Ga$ fibre
into the $\Ga^{d_0}$ factor gives a presentation with $r=0$, so case~(c) for $r\le1$ needs only the onto
case of~(a). For $d_0=0$, $r\le2$ the engine runs directly on~(b).
\end{itemize}

\subsection{The bridges to \texorpdfstring{\texttt{lean\_mstar}}{lean\_mstar}}
From Theorem~\ref{thm:vec} we prove, unconditionally, \lean{lean_mstar}'s six vector-extension hypotheses
(\lean{LSG/W88/V5dOntoBridges.lean}, \lean{V5dSMGn.lean}, \lean{V5dMXPMain.lean}): \lean{Thm11KappaMXG}
(Theorem~MX), \lean{Thm11KappaMGn} and \lean{Thm11KappaMGnK} for every CM lattice, \lean{Thm11KappaNat}
(the four exponentials statement at $t=c\vp$), \lean{Thm11KappaSMGn} (Theorem~S), and \lean{Thm11KappaMXP}
(Theorem~MX$'$, which has $r=2$ and uses case~(b)). As in \lean{lean_mstar}, the statements for a general CM
lattice carry the premise $\kappa_1=0$ (\lean{KappaOneZero}), proved there for both lattices used. The CM
inputs are proved in Lean: rotation formulas, quasi-periods via Legendre, \lean{LiftAlg} and
\lean{PerTanRat} for CM blocks, and the splitting of abelian subvarieties of $E_i^n\times E_\rho^n$ via the
irrationality of $\sqrt3$. The statements are verbatim copies of \lean{lean_mstar}'s
(\lean{LSG/W88/MStarCopy*.lean}); \lean{lean_mstar}, which has this project as a Lake dependency, derives its
own unmodified hypotheses from them, so a drift of a copy breaks its build. Two of the CM abelian-variety
hypotheses of \lean{lean_mstar} are treated in Section~\ref{sec:abvar}.

\subsection{Fidelity review}
A second AI session that had not written the proofs reviewed the vector-extension statements against
\cite[pp.~375, 385--389, 394--395]{W88}. It found no fatal or major issue: $\Lambda(G)$, $\Omega(G)$, the
$\Qbar$-structure of the push-out and the subgroup family match Waldschmidt's for this class (including
completeness of the family), and all casts are safe. As a test of non-vacuity it derived in Lean, from each
of (a), (b) and~(c), that $\eta(\omega_1)\notin\Qbar\,\omega_1$ for the lemniscatic curve (with $G=E^\natural$). Its main
remark is the scope of the claim, stated above; the others concern docstrings, the redundancy of
\lean{Proper}, and the fact that the lattice case of \lean{LamV} depends on the junk value of $\zeta$ at
lattice points (proved in the same file, so a change of definition would break the build).

% ===========================================================================
\section{An abelian surface: the Jacobian of \texorpdfstring{$y^5=x(1-x)$}{y\^{}5=x(1-x)}}\label{sec:abvar}
% ===========================================================================

\subsection{The group and the statements}
Theorem~A of \cite{MstarPaper} at the Beta value $B(1/5,1/5)$ applies Waldschmidt's theorem to
$G=\Ga^{d_0}\times\Gm^{d_1}\times A^r$, where $A$ is the Jacobian of the genus-two curve
$C\colon y^5=x(1-x)$, a simple abelian surface with complex multiplication by $\QQ(\zeta_5)$. Here $A$ is the
Jacobian of $C$, presented in Lean by its period lattice in the coordinates $\int dX/Y$, $\int X\,dX/Y$ on the
model $Y^2=4X^5+1$ ($X=-y$, $Y=2x-1$). Lean computes these periods (Beta values) and proves that they are
periods of an analytic exponential map onto Grant's model in $\PP^8$ with discrete kernel. The
identification of that torus with $\operatorname{Jac}(C)$, and of complex subtori with abelian subvarieties,
is classical and not formalized. The final theorems (e.g.\ \lean{two_three_beta5_uncond} below) are statements
about explicit numbers and do not depend on it.

Concretely, $T_G=\CC^{d_0}\times\CC^{d_1}\times(\CC^2)^r$, and \lean{latt}${}\subset\CC^2$ is the lattice spanned by
the four periods $\int_{\gamma_e}(dX/Y,\,X\,dX/Y)=(c_0\zeta^e,\,c_1\zeta^{2e})$, $e=0,\dots,3$, where $\zeta=e^{2\pi i/5}$,
$c_j=\frac{(-1)^j}5(1-\zeta^{j+1})B\bigl(\frac{j+1}5,\frac{j+1}5\bigr)$ and $\gamma_e$ are explicit loops. Lean proves
that the four periods are $\RR$-independent, that \lean{latt} lies in the kernel of the exponential map, and that
this kernel lies in $\QQ\cdot$\lean{latt}; so \lean{latt} is a finite-index sublattice of the period lattice, and
every statement below depends only on its $\QQ$-span. The data are restricted as follows:
\[
\Lambda_A=\{(a,b,z): a\in\Qbar^{d_0},\ e^{b_i}\in\Qbar,\ z_k\in\texttt{latt}\},\qquad
\Omega_A=0\times2\pi i\ZZ^{d_1}\times\texttt{latt}^r ,
\]
and a subspace $V\subset T_G$ is \emph{toric-injective} if every $v\in V\cap\QQ\Omega_A$ with zero toric part is
$0$, i.e.\ $V$ contains no nonzero $\QQ$-period of $A^r$ alone. Connected algebraic subgroups are products
$H_0\times T_U\times B$, with the tangent space $t$ of $B$ characterized by ``$t\cap\texttt{latt}^r$ spans $t$
over $\RR$''. The statements are in \lean{LSG/W88/StatementAbVar.lean}.

\begin{theorem}[Lean; standard axioms]\label{thm:abvar}
Let $G=\Ga^{d_0}\times\Gm^{d_1}\times A^r$, $A$ as above, so $d_2=2r$.
\begin{enumerate}
\item[(a)] Theorem~\ref{thm:W88} holds under \NR{} for toric-injective $V$ and $Y\subset V\cap\Lambda_A$
(\lean{LSG.W88.theorem41AbVarPer}).
\item[(b)] The \S5d consequence ((5.4) with $d_2=2r$, without \NR) holds for toric-injective hyperplanes $V$
and $y_1,\dots,y_m\in V\cap\Lambda_A$ (\lean{LSG.W88.hyperplaneKappaAbVar}); by a common-denominator argument
the $y_i$ may be taken in $V\cap\QQ\Lambda_A$.
\end{enumerate}
\end{theorem}

\emph{Scope.} Theorem~\ref{thm:abvar} is W88 \S5d (and Thm.~4.1 under \NR) for $\Ga^{d_0}\times\Gm^{d_1}\times A^r$,
$A=\operatorname{Jac}(y^5=x(1-x))$, for toric-injective hyperplanes $V$ and subgroups $Y$ whose $A$-components
are (rational) periods; so $\exp_G(Y)\subset\Ga^{d_0}(\Qbar)\times\Gm^{d_1}(\Qbar)\times A^r_{\mathrm{tors}}$. Both
restrictions only weaken Waldschmidt's theorem, and they are all that \lean{lean_mstar}'s models use. For
$G=A^r$ alone ($d_0=d_1=0$) the restricted statement is empty, since toric injectivity forces
$V\cap\Lambda_A=0$; it must not be cited as Waldschmidt's theorem for $A^r$.

\subsection{The proof}
The architecture is again that of Sections~\ref{sec:stage2}--\ref{sec:5d}, but every layer that touched the
elliptic curve had to be redone for a surface without a Weierstrass-type uniformization in Mathlib.
\begin{itemize}
\item \emph{An explicit model and certificates} (\lean{LSG/Jac/}). We use Grant's model $V(Q)\subset\PP^8$ of the
Jacobian of a genus-two curve with a rational Weierstrass point (quadrics in Mumford coordinates), with a
complete system of addition laws of bidegree $(2,2)$, the duplication map and the invariant vector fields.
The polynomial identities behind these facts (laws preserve the model, flow invariance, the doubling
identities, a degree-$7$ Nullstellensatz showing that duplication does not vanish on $V(Q)\setminus0$, local
complete-intersection systems) are far beyond \texttt{ring}. They are produced by Python scripts with exact
arithmetic, which are untrusted, and checked in Lean by a normal-form checker modulo the ideal of the
model, with Kronecker-packed natural-number accumulators. Each check is a Boolean computation proved by
\texttt{decide\ +kernel}, so the kernel evaluates it, and a soundness theorem turns a successful check into the
polynomial identity over any commutative ring. There is no \texttt{native\_decide}. The certificates take
about five CPU-hours to check.
\item \emph{The analytic exponential} (\lean{LSG/JacAn*/}). We build the genus-two Abel map and a local
exponential near a generic pair of points, prove Abel's theorem locally (Euler--Jacobi), extend the
exponential to an entire map $\CC^2\to\PP^8$ by doubling the germ, prove that it is a homomorphism, that
its kernel is discrete and contains \lean{latt} (via path continuation of Abel integrals, including through
Weierstrass points), and compute the Beta periods. An entire nowhere-vanishing lift $\theta$, normalized
to have an affine factor of automorphy and order two, plays the role of $\sigma$; a Liouville argument shows
that the addition laws are correct and complete on its image, and a boundary flow-exit argument proves that
every nonzero point of the model is $c\,\theta(z)$. Taylor coefficients at algebraic points are algebraic,
with explicit size bounds, in three charts that cover the model.
\item \emph{A generic multiplicity estimate} (\lean{LSG/MEB/}). The estimate of Section~\ref{sec:stage2} is
proved once more for $\Ga^{d_0}\times\Gm^{d_1}\times B_1\times\dots\times B_m$, generic over an interface of
\emph{abelian blocks}: an entire lift with discrete period group, an explicit ideal whose zero set is the
cone over the image, addition laws correct on the image, chart flows, and Hilbert and local-CI
certificates. The interface is semantic: identities that the elliptic proof checked formally modulo the
cone become consequences of the laws on the image, which avoids an estimated twenty- to forty-fold increase
of the certificate work. The Jacobian is the instance \lean{LSG.JacBlock.aBlock}.
\item \emph{The analytic side and \S7} (\lean{LSG/AuxFunA/}, \lean{LSG/W88/AbVar*}). For a simple abelian
surface, W83's Prop.~2.1 in general needs a several-variable interpolation estimate of Masser type. Toric
injectivity avoids it: $\ZZ$-independent periods in a toric-injective $V$ have $\QQ$-independent toric parts,
hence are $\CC$-independent, and the linear form of Prop.~2.1 used for elliptic curves applies with
$\theta$'s factor of automorphy as multiplier. Lemma~U needs Taylor data only at the origin, because every
point of $\Lambda_A$ has its $A$-part in \lean{latt}. W88 \S7 and the \S5d engine are written generically
for abelian blocks (\lean{theorem41B_of}); the \S5d maximizer must be defined over $\Qbar$, and we prove that
abelian tangent spaces of $A^r$ are $\Qbar$-rational using the CM structure (\lean{ratTangentsA}).
\end{itemize}

\subsection{The bridge to \texorpdfstring{\texttt{lean\_mstar}}{lean\_mstar}}
From Theorem~\ref{thm:abvar}(b) we prove \lean{lean_mstar}'s hypotheses \lean{AbVar.Thm11KappaA} and
\lean{AbVar.Thm11KappaAK} at $y^5=x(1-x)$ (\lean{LSG.W88.thm11KappaA5}, \lean{thm11KappaAK5}), stated against
verbatim copies. The transport rescales \lean{lean_mstar}'s coordinates, which are dual to $dx/y^{4-s}$, by
the algebraic factors $5(-1)^s$, and checks that \lean{lean_mstar}'s toric injectivity implies ours and that
every nonzero abelian tangent lands in \lean{lean_mstar}'s subgroup family. In \lean{lean_mstar} (which has this
project as a Lake dependency) these yield, with no hypothesis and standard axioms only:
\lean{two_three_beta5_uncond}, that $2^{t}$ and $3^{t}$ are not both algebraic for $t=B(1/5,1/5)/\pi$; the
four exponentials statements at $cB(1/5,1/5)/\pi$ and $c\pi/B(1/5,1/5)$ for algebraic $c\ne0$; and the input
$(1-\zeta_5)B(1/5,1/5)/\pi\notin\Qbar$ (Wüstholz), here derived from Waldschmidt's theorem.

\subsection{Fidelity review}
A further AI session that had not written the proofs reviewed the statements and the bridge against
\cite[pp.~375, 385--386, 394--395]{W88} and \lean{lean_mstar}. It found no fatal or major issue. Its main remark is
the scope above (the second restriction, on the $A$-components, had been omitted from a sentence in
\lean{lean_mstar}'s paper); it also verified in Lean that the subgroup condition is not vacuous (a nonzero
subspace of $\CC^2$ that is not an abelian tangent) and that the restricted statement is empty for $A^r$
alone, and asked for the wording of ``$A$ is the Jacobian'' used above.

% ===========================================================================
\section{A vector extension of the Jacobian}\label{sec:abvarv}
% ===========================================================================

\subsection{The group and the statements}
Theorem~MG$^\natural$-A of \cite{MstarPaper} at $B(1/5,1/5)$ gives the headline statement that $2^t$ and $3^t$ are
not both algebraic for $t=B(1/5,1/5)=\Gamma(1/5)^2/\Gamma(2/5)$ (without the division by $\pi$ of
Section~\ref{sec:abvar}). It applies Waldschmidt's theorem to a group with a vector part. Let
$A=\operatorname{Jac}(y^5=x(1-x))$ as in Section~\ref{sec:abvar}, and let $A_{\bar\sigma}^\natural$ be the extension of $A$
by $\Ga$ whose class is that of the differential of the second kind $dx/y$; on the model $Y^2=4X^5+1$ this is the
class $[X^3\,dX/Y]=-\tfrac15[dx/y]$, the $\bar\sigma$-eigenline for the action of $\QQ(\zeta_5)$. Let $Q$ be the push-out
of $(A_{\bar\sigma}^\natural)^m$ along $b^{\mathsf T}$, for $b$ an adapted basis of a number field containing $\QQ(\zeta_5)$
(that is, the numbers $\zeta^e b_j$ are $\QQ$-linearly independent). Then $T_Q=(\CC^2)^m\times\CC$, and the periods of $Q$
are the vectors $(\lambda,\sum_j b_j\bar\eta(\lambda_j))$, with $\lambda\in$ \lean{latt}${}^m$ and $\bar\eta$ the fibre
period.

The statements are in \lean{LSG/W88/AbVarVStmt.lean}. They are generic over a family of \emph{vector blocks}
$B_1,\dots,B_m$ (an abelian block of Section~\ref{sec:abvar} together with fibre data) and a matrix $M\colon\CC^m\to\CC^r$.
The data are restricted as in Section~\ref{sec:abvar}: $\Lambda_{A^\natural}$ consists of the $(a,b,z,s)$ with
$a\in\Qbar^{d_0}$, $e^{b_i}\in\Qbar$, $z_k\in$ \lean{latt} and $s-M\bar\eta(z)\in\Qbar^r$, so the $A^m$-component of every
point is $0$ while its fibre component is an arbitrary algebraic point; and $V$ is \emph{toric-injective} if it
contains no nonzero rational period of $Q$.

\begin{theorem}[Lean; standard axioms]\label{thm:abvarv}
Let $G=\Ga^{d_0}\times\Gm^{d_1}\times Q$ with $Q$ as above, so $d_2=2m+1$.
\begin{enumerate}
\item[(a)] The \S5d consequence ((5.4), without \NR) holds for toric-injective hyperplanes $V$ and
$y_1,\dots,y_\ell\in V\cap\Lambda_{A^\natural}$ (\lean{LSG.W88.hyperplaneKappaAbVarV_aBlockV}).
\item[(b)] More generally, Theorem~\ref{thm:W88} holds under \NR{} for $\Ga^{d_0}\times\Gm^{d_1}\times(B_1^\natural\times\dots
\times B_m^\natural\text{ pushed out along }M)$, for every family of vector blocks with analytic and Taylor
packages, every onto algebraic $M$ with $r\le1$, toric-injective $V$, and $Y$ whose base parts lie in the period
lattice (\lean{LSG.W88.theorem41AbVarVPer}).
\item[(c)] Philippon's multiplicity estimate holds for these groups, generic over vector blocks
(\lean{LSG.MEBV.multEstV_holds}).
\end{enumerate}
\end{theorem}

\emph{Scope.} For Theorem~MG$^\natural$-A at $B(1/5,1/5)$ the formal development proves the following special case of
W88 \S5d: Theorem~\ref{thm:abvarv}(a), that is, W88 \S5d for $G=\Ga^{d_0}\times\Gm^{d_1}\times Q$, restricted to
\begin{itemize}
\item hyperplanes $V$ that are \emph{toric-injective} ($V$ contains no nonzero rational period of $Q$);
\item subgroups $Y$ whose points have trivial $A^m$-component (their $\Ga$, $\Gm$ and fibre components are arbitrary
algebraic points).
\end{itemize}
This is implied by W88. The push-out is one-dimensional ($r=1$) and $b$ is adapted; this is what guarantees that
every algebraic subgroup with nonzero abelian part contains the fibre. The generic Props behind
Theorem~\ref{thm:abvarv} are instances of W88 only when $M$ is algebraic and, for (a), when the block family has this
fibre property; the proved cases satisfy both. Unlike Section~\ref{sec:abvar}, the restricted statement is not empty
for $d_0=d_1=0$, since non-torsion points of the fibre are allowed. The Lean proof, which goes through Philippon's
multiplicity estimate, is unconditional. The hypothesis of Theorem~MG$^\natural$-A is used only at three frame
hyperplanes, and each of them is toric-injective. This is proved in Lean from the $F$-independence of $b$ and the
fact that the fibre character $\zeta\mapsto\zeta^4$ is a field embedding.

The group $A_{\bar\sigma}^\natural$ is presented in Lean by an explicit projective model over $\QQ$ in $\PP^{11}$
(Grant's variety, three fibre coordinates $T_0,T_4,T_8$, Plücker quadrics and $81$ addition laws of bidegree $(3,3)$)
and an entire uniformisation $(z,w)\mapsto(\theta_N(z),T(z,w))$ built from theta functions. Its fibre periods are
computed in Lean to be $\bar\eta([\zeta^e]\gamma)=\frac35\zeta^{4e}(1-\zeta^4)B(4/5,4/5)$, i.e.\ $-3$ times the periods of
$X^3\,dX/Y$. The identification of this model with the extension of $\operatorname{Jac}(C)$ is classical and not
formalised. The final four-exponentials statement concerns explicit numbers and does not depend on it.

\subsection{The proof}
The architecture is that of Section~\ref{sec:abvar}, with a fibre added to every layer.
\begin{itemize}
\item \emph{The analytic group} (\lean{LSG/JacAnV/}). A Liouville argument gives
$\partial_i\partial_j\log\theta_{N,0}=-3\wp_{ij}+c_{ij}$ from six small quadric identities. The fibre coordinates are
the derivative lifts $T_k(z,w)=\theta_k(z)(w+\ell(z))-\frac13\partial_0\theta_k(z)$. They have the same factor of
automorphy as $\theta_N$, with the fibre period $\bar\eta$ as translation in $w$, and their flows are given by rational
quadrics. The fibre period formula comes from the path-continuation machinery of Section~\ref{sec:abvar}.
\item \emph{The model} (\lean{LSG/JacBlockV/}, \lean{LSG/Jac/B6V*}). The fibre addition laws are the polarization of
the nine certified laws of $A$, defined with \lean{pderiv} and proved correct by differentiating them, so no new law
certificates are needed. The only new certificates are those of the Hilbert bound: every homogeneous polynomial of
degree $41$ lies in the ideal of the model plus four linear forms, from eight kernel-checked identities.
\item \emph{A generic multiplicity estimate} (\lean{LSG/MEBV/}). The estimate is proved once more, generic over an
interface \lean{VBlock} of vector blocks. Its design builds in five corrections found by an exact computer-algebra
spike: slicing in the $X$-variables only (the model is not a local complete intersection on $\{X=0\}$), a Hilbert
certificate with linear forms, fibre flows only in their own chart, generators that need not span a saturated
ideal, and no requirement on the fibre coordinate of the identity.
\item \emph{The analytic side, \S7 and \S5d} (\lean{LSG/AuxFunAV/}, \lean{LSG/W88/AbVarV*}). Under toric injectivity,
W83's Prop.~2.1 again needs only its linear form, and the fibre forms are quasi-periodic with $\theta$'s multiplier.
The \S5d engine needs three properties of the block family, all proved for $A_{\bar\sigma}^\natural$ from the CM
structure: the fibre property above, an algebraic relation between the fibre periods on a rational abelian tangent
and on its pivot coordinates, and $\Qbar$-rationality of abelian tangents.
\end{itemize}

\subsection{The bridge to \texorpdfstring{\texttt{lean\_mstar}}{lean\_mstar}}
From Theorem~\ref{thm:abvarv}(a) we prove \lean{lean_mstar}'s hypotheses \lean{AbVar.Thm11KappaAn} and
\lean{AbVar.Thm11KappaAnK} (the latter for $L_{A^\natural}=\emptyset$) at $y^5=x(1-x)$
(\lean{LSG.W88.thm11KappaAn5}, \lean{thm11KappaAnK5}). Like \lean{Thm11KappaA}, these hypotheses are restricted to
toric-injective hyperplanes, a restriction \lean{lean_mstar} adopted for this purpose (its frame hyperplanes are
toric-injective, \lean{perAn_inj}). The transport rescales by algebraic factors and maps our fibre periods onto
\lean{lean_mstar}'s generators, since $\frac53\cdot b_j\frac35\zeta^{4e}\bar P=b_j\zeta^{(5-1)e}\bar P$. In
\lean{lean_mstar}, \lean{Bridge.two_three_beta5'_uncond} (that $2^t$ and $3^t$ are not both algebraic for $t=B(1/5,1/5)$) and
\lean{Bridge.wuestholzB5_uncond} ($(1-\zeta^4)B(4/5,4/5)/\pi\notin\Qbar$) follow with standard axioms only.

\subsection{Fidelity review}
A further AI session that had not written the proofs reviewed the statements, the block data and the bridge against
\cite[pp.~385--386, 394--395]{W88} and \lean{lean_mstar}. It found no fatal or major issue. Its main remarks are the
scope above and that \lean{lean_mstar} had yet to adopt the restriction, which it has since done. It proved in Lean, in a scratch file, that
\lean{lean_mstar}'s frame hyperplanes are toric-injective, and it re-ran the three \lean{lean_mstar} consumers under the
narrowed hypothesis. It also checked non-vacuity: the fibre period $\bar\eta$ is nonzero, the data contain non-torsion
fibre points, and the subgroup family of the model is exactly the family of connected algebraic subgroups.

% ===========================================================================
\section{Explicit consequences, formally verified}\label{sec:conseq}
% ===========================================================================

The certificate \lean{lean_mstar} \cite{LeanMstar}, which has this project as a Lake dependency, derives explicit
transcendence statements from the hyperplane theorems of Sections~\ref{sec:5d}, \ref{sec:vecext}, \ref{sec:abvar} and
\ref{sec:abvarv}. We list them with their status at its commit \lean{7ddd73d}. Lean names are in the namespace \lean{CM4E}
unless they start with \lean{LSG}; each entry names the sub-namespace once.

\emph{Notation.} $\mathcal L=\{z\in\CC:e^z\in\Qbar^\times\}$, $\widetilde{\mathcal L}=\Qbar+\Qbar\mathcal L$, and
$\mathcal L_E$ is the set of elliptic logarithms of algebraic points of $E$. The curves are $E_i$ (lattice $\vp\ZZ[i]$) and
$E_\rho$ (lattice $\Omega_\rho\ZZ[\rho]$, $\Omega_\rho=\Gamma(1/3)^3/(2^{4/3}\pi)$), and $\eta_\rho=\eta(\Omega_\rho)$.
$\mathrm{4E}(t)$ (\lean{FourExpAt}) says: if $a,b\in\mathcal L$ are $\QQ$-linearly independent, then $e^{ta},e^{tb}$ are not
both algebraic. $A$ is the Jacobian of $y^5=x(1-x)$ and $B=B(1/5,1/5)=\Gamma(1/5)^2/\Gamma(2/5)$. Throughout, $c\in\Qbar^\times$.

\emph{Derivation.} Each statement comes from a group $G$, a $\Qbar$-rational map $\Psi\colon T_G\to\CC^2$ and a line $V$ of
transcendental slope. For a number field $K$ of degree $n$, $G''$ is the quotient of $G^n$ by the largest connected algebraic
subgroup with tangent space in the kernel of $\Psi\circ(z_j)_j\mapsto\sum_j\theta_jz_j$, for a basis $\theta$ of $K/\QQ$. With
$\alpha(H)=2\dim H-2d_0(H)-d_1(H)$, $\alpha_{K,\Psi}=\alpha(G'')/n$ and $\kappa=\dim_K(V\cap\Psi(\Omega_K(G)))$, the bound is
$\dim_K(V\cap\Psi(\Lambda_K(G)))\le\alpha_{K,\Psi}-\kappa$ for the $K$-spans $\Lambda_K$, $\Omega_K$, and a counterexample
would give more $K$-independent data than that. Table~\ref{tab:conseq} has the numbers for large $n$, as computed in
\cite{MstarPaper}. The Lean proofs do not use this bound: they apply (5.4) of \cite[\S5d]{W88} to a hyperplane of a model of
$G''$.

\begin{table}[ht]\footnotesize
\begin{tabular}{@{}lllllll@{}}
\toprule
Entry & $G''$ & $\alpha_{K,\Psi}$ & $\kappa$ & data & bound & Status \\
\midrule
MG & $\Gm^{2n}\times A''$, $A''=E^n/U$ & $3$ & $1$ & $3$ & $2$ & \textbf{U}, \textbf{I} at $\Omega_\rho$ \\
MG$^\natural$ & $\Gm^{2n}\times Q''$, $Q''=(E^\natural)^n/U$ & $3+2/n$ & $1$ & $3$ & $2+2/n$ & \textbf{U}, \textbf{I} at $\Omega_\rho$ \\
SMG(a) & $\Ga^2\times\Gm^{2n}\times A''$ & $3$ & $1$ & $3$ & $2$ & \textbf{U}$^*$ \\
SE(1) & $\Ga^2\times A_1''\times A_2''\times A'''$ & $3$ & $1$ & $3$ & $2$ & \textbf{I} \\
MX, deg2(ii) & $(E_i^\natural\times E_\rho^\natural\times\Ga^2)^n/U$ & $4+2/n$ & $2$ & $3$ & $2+2/n$ & \textbf{I} \\
MX$'$, deg2(i) & ext.~of $E_i^n\times E_\rho^n$ by $\Ga^2$ & $4+4/n$ & $2$ & $3$ & $2+4/n$ & \textbf{I} \\
Theorem M & $A''\times\Gm^n\times\Ga$ & $2$ & $1$ & $2$ & $1$ & \textbf{U}, \textbf{I} \\
MG-A at $B/\pi$ & $\Gm^{2n}\times A''$, $A''\sim A^r$ & $3$ & $1$ & $3$ & $2$ & \textbf{U} \\
MG$^\natural$-A at $B$ & $\Gm^{2n}\times Q''$, $Q''=(A^\natural)^r/U$ & $3+2/n$ & $1$ & $3$ & $2+2/n$ & \textbf{U} \\
X7 & $\prod_jA_j''\times\Ga^2$ & $4$ & $2$ & $3$ & $2$ & \textbf{W} \\
X9 & $\prod_jA_j''\times\Ga^2$ & $4$ & $2$ & $3$ & $2$ & \textbf{N} \\
X$^\kappa$ & $\prod_jA_j''\times\Ga^2$ & $2\kappa$ & $\kappa$ & $\kappa+1$ & $\kappa$ & \textbf{N} \\
\bottomrule
\end{tabular}
\caption{The consequences. ``Data'' counts the period data and the vectors a counterexample would give.}\label{tab:conseq}
\end{table}

\emph{Status labels.} \textbf{U}: verified with no hypothesis, that is, a declaration whose only premises belong to the
mathematical statement. \textbf{U}$^*$: every premise of a generic declaration is discharged at this commit, but no single
declaration states the composite. \textbf{I}: verified with no form of W88 assumed, modulo the named classical inputs listed,
which are explicit hypotheses. \textbf{W}: verified conditionally on a named form of W88 that this project does not prove.
\textbf{N}: not formalised as
stated. Two inputs recur: $(g_3)$ is $g_3(\Omega_\rho\ZZ[\rho])\in\Qbar$ (the argument \texttt{hg$_3'$}; $g_2=0$ is proved, and
$g_3=4$ is classical but not formalised), and $(\Gamma)$ is $\vp\Omega_\rho/\pi\notin\Qbar$ (\texttt{h$\Gamma$};
W\"ustholz \cite{Wustholz1989}).

\begin{itemize}
\item \textbf{MG.} $\mathrm{4E}(c\vp/\pi)$; for example $2^{\vp/\pi},3^{\vp/\pi}$ are not both algebraic. For $s=c\vp/\pi$ and
$K\ni i,c$, no $x_1,x_2\in K\mathcal L$ have $sx_j\in K\mathcal L+K\mathcal L_{E_i}$ with $2\pi i,x_1,x_2$ $K$-independent.
\lean{Unconditional}: \lean{fourExpAt_varpi_div_pi_all}, \lean{two_three}, \lean{mg_K_varpi}, from \lean{LSG.W88.thm11Kappa},
\lean{LSG.W88.hypStdCMAt}: \textbf{U}. At $\Omega_\rho$, \lean{fourExpAt_omegaRho_div_pi_all}: \textbf{I} ($(g_3)$). Arbitrary
$K\supseteq\QQ(i)$: \textbf{N}.
\item \textbf{MG$^\natural$.} $\mathrm{4E}(c\vp)$; for example $2^{\vp},3^{\vp}$ are not both algebraic. \lean{Unconditional}:
\lean{fourExpAt_varpi_all}, \lean{two_three_varpi}, the $K$-form \lean{mgn_K_varpi}, from \lean{LSG.W88.V5d.thm11KappaMGn_uncond},
\lean{thm11KappaNat_uncond}, \lean{thm11KappaMGnK_uncond}: \textbf{U}. At $\Omega_\rho$, \lean{fourExpAt_omegaRho_all},
\lean{two_three_omegaRho}: \textbf{I} ($(g_3)$).
\item \textbf{SMG(a).} For $s=c\vp/\pi$ and $y_1,y_2\in\widetilde{\mathcal L}$ with $2\pi i,y_1,y_2$ $\Qbar$-independent, $sy_1,sy_2$
are not both in $\widetilde{\mathcal L}$. \lean{Gen.smg_false} at \lean{Gen.latI}, with \lean{Bridge.thm11KappaSMG_of_std} and
\lean{LSG.W88.hypStdCMAt} ($d_0=2$): \textbf{U}$^*$; the composite occurs inside \lean{Gen.thmS}, hence in
\lean{Unconditional.thmS_varpi}. With second row in $\widetilde{\mathcal L}+\Qbar\mathcal L_{E_i}$: \textbf{N}.
\item \textbf{SE(1).} For $t=c\Omega_\rho/\vp$ and $u_1,u_2\in\Qbar+\Qbar\mathcal L_{E_i}$ with $\vp,u_1,u_2$ $\Qbar$-independent,
$tu_1,tu_2$ are not both in $\Qbar+\Qbar\mathcal L_{E_i}+\Qbar\mathcal L_{E_\rho}$. \lean{Unconditional.thmSE_varpi_omegaRho}, from
\lean{LSG.W88.hypStd_of_pair}: \textbf{I} ($(g_3)$). For two non-isogenous CM lattices, \lean{Bridge.thmSE_of_nonIsog} assumes no
form of W88 but has the premise \lean{Bridge.CMMulLogE}, proved only for $E_i$, $E_\rho$.
\item \textbf{MX, MX$'$, deg2.} Let $X\in\Lambda(E_i^\natural)$, $Y\in\Lambda(E_\rho^\natural)$ with base and fibre coordinates $(z,w)$,
$a,a'\in\Qbar$, $\beta\in\Qbar^\times$. (MX) $(z(X)+z(Y)+a)\beta\eta_\rho=(w(X)+w(Y)+a')\vp$ implies
$z(X)+z(Y)+a\in\Qbar\vp+\Qbar\Omega_\rho$, and then $a=0$, $z(X)\in\QQ(i)\vp$, $z(Y)\in\QQ(\rho)\Omega_\rho$. (MX$'$)
$(z(X)+w(Y)+a)\beta\Omega_\rho=(w(X)+z(Y)+a')\vp$ implies $z(X)+w(Y)+a\in\Qbar\vp+\Qbar\eta_\rho$ (Proposition~D$'$ of
\cite{MstarPaper}; the step to its Theorem~MX$'$ is \textbf{N}). (deg2)(i) $\pi,\vp,\Omega_\rho,\vp\Omega_\rho$ are
$\Qbar$-linearly independent; (ii) so are $1,\vp,\Omega_\rho,\vp\Omega_\rho/\pi$. \lean{Unconditional}, all \textbf{I}, all with
$(g_3)$, $(\Gamma)$: \lean{mx_K}; \lean{mx} (also \lean{Gen.WuestholzMX}, \lean{Gen.MasserCM} \cite{Masser1975}); \lean{mx1}, its
corollary on $\zeta$ (also \lean{Gen.LaurentZeta} \cite{Laurent1980}); \lean{mxP_K}; \lean{deg2_i}, \lean{deg2_ii} (also
\lean{Gen.WuestholzSix}, which implies $(\Gamma)$). From \lean{LSG.W88.V5d.thm11KappaMXG_uncond} ($r=1$) and
\lean{thm11KappaMXP_uncond} ($r=2$). These statements use hand-written quasi-periods and the coordinates \lean{MX.LogNat},
checked numerically in \lean{lean_mstar}'s reviews.
\item \textbf{Theorem M.} If $u\in\Qbar\mathcal L_{E_i}$, $\ell\in\widetilde{\mathcal L}$ and $i\pi u=c\vp\ell$, then $u=sc\vp$,
$\ell=si\pi$ with $s\in\Qbar$; hence $\wp(c\vp/\pi)\notin\Qbar$ (M1), and $\wp(c\pi\vp)\notin\Qbar$ if $e^{\pi^2}\in\Qbar$ (M3).
\lean{Unconditional}: \lean{thmM_varpi}, \lean{m1_varpi}, \lean{m3_varpi}, from \lean{LSG.W88.hypStdCMAt} ($d_0=1$): \textbf{U}.
Also \textbf{I}: \lean{thmM_varpi_k} ($s\in\QQ(i)$ for $u\in\mathcal L_{E_i}$, $c=1$), \lean{meta_varpi} (\lean{Gen.MasserCM});
\lean{m3'_varpi} (\lean{Gen.HermiteLindemann}); \lean{thmM_omegaRho} ($(g_3)$). (M2) of \cite{MstarPaper}: \textbf{N}.
\item \textbf{MG-A at $B/\pi$.} $\mathrm{4E}(cB/\pi)$ and $\mathrm{4E}(c\pi/B)$; for example $2^t,3^t$ are not both algebraic at
$t=\Gamma(1/5)^2/(\pi\Gamma(2/5))$. Also $P/\pi\notin\Qbar$ for $P=(1-\zeta_5)B$, and the $K$-form at $P$ with $z_\sigma(\Lambda(A))$
replaced by $\ZZ P$. \lean{Bridge}: \lean{two_three_beta5_uncond}, \lean{fourExpAt_beta5_uncond},
\lean{fourExpAt_pi_div_beta5_uncond}, \lean{wuestholzA5_uncond}, \lean{mgA5_K_uncond}, from \lean{LSG.W88.thm11KappaA5},
\lean{LSG.W88.thm11KappaAK5}, within the scope of Theorem~\ref{thm:abvar}: \textbf{U}. One of $e^{a\pi^2},e^{b\pi B},e^{c\pi^2/B}$
is transcendental ($a,b\in\Qbar^\times$), \lean{Bridge.corA2_beta5_uncond}: \textbf{I} (\lean{AbVar.NotLtildeA}, that is,
$P\notin\widetilde{\mathcal L}$). The $K$-form with logarithms of $A$ among the data, \lean{AbVar.mgA_K}: \textbf{W}
(\lean{AbVar.Thm11KappaAK}).
\item \textbf{MG$^\natural$-A at $B$.} $\mathrm{4E}(cB)$ and $\mathrm{4E}(c/B)$; for example $2^t,3^t$ are not both algebraic at
$t=B=\Gamma(1/5)^2/\Gamma(2/5)$. Also $\bar P/\pi\notin\Qbar$ for $\bar P=(1-\zeta_5^4)B(4/5,4/5)$, and the $K$-form at $\bar P$ with the
fibre values replaced by $\ZZ\bar P$. \lean{Bridge}: \lean{two_three_beta5'_uncond}, \lean{fourExpAt_beta5'_uncond},
\lean{fourExpAt_div_beta5_uncond}, \lean{wuestholzB5_uncond}, \lean{mgAn5_K_uncond}, from \lean{LSG.W88.thm11KappaAn5},
\lean{LSG.W88.thm11KappaAnK5}, within the scope of Theorem~\ref{thm:abvarv}: \textbf{U}. The $K$-form with logarithms of
$A^\natural$ among the data, \lean{AbVar.mgAn_K}: \textbf{W} (\lean{AbVar.Thm11KappaAnK}).
\item \textbf{X7 and genus three.} On the Jacobian of $y^7=x(1-x)$, with $P_j=(1-\zeta_7^j)B(j/7,j/7)$ and $E$ with CM by
$\QQ(\sqrt{-7})$ and period $\Omega$: $(z_1(v_1)+u+q_1)\beta P_3=(z_3(v_3)+z_2(v_2)+q_2)P_1$ implies $q_1=q_2=0$,
$z_j(v_j)\in\Qbar P_j$, $u\in\Qbar\Omega$, and all points are torsion. Also $1,B(\frac27,\frac27),B(\frac37,\frac37),
B(\frac37,\frac37)/B(\frac17,\frac17)$ are $\Qbar$-linearly independent. \lean{AbVar}: \lean{thmX7}, \lean{thmX7_torsion},
\lean{x7c}, \lean{x7a_AJ}, \lean{x7b_AJ}, with logarithms as a parameter set \lean{LA}: \textbf{W}, with hypotheses
\lean{Thm11KappaX7}, \lean{ChowlaSelberg7}, \lean{WuestholzBeta7}; except for \lean{x7c} also \lean{WuestholzX7}; for the torsion
clause and the Abel--Jacobi forms \lean{x7a_AJ}, \lean{x7b_AJ} also \lean{TorsionX7}. Also \textbf{W}:
$\mathrm{4E}(cB(\frac17,\frac17)/\pi)$ (\lean{fourExpAt_beta7}, hypothesis \lean{Thm11KappaA} at $N=7$). This project does not cover
genus three: it proves no form of these hypotheses, and \lean{dev/M3.md} is a design note with no Lean behind it.
\item \textbf{X9.} The level-$9$ analogue of X7, on the CM threefold factor of the Jacobian of $y^9=x(1-x)$ and $E_\rho$
\cite{MstarPaper}.
\textbf{N}: the level-$9$ data are not in Lean. The generic \lean{AbVar.XSetup.thmXA} (no torsion clause, the period relation
as a premise) is \textbf{W} (\lean{AbVar.Thm11KappaXA}, \lean{AbVar.WuestholzXA}, \lean{AbVar.WuestholzSplitXA}); this project
proves none of these.
\item \textbf{X$^\kappa$.} The statement with $\kappa=2^{g-2}$ parallel pairs of slots at level $N$ \cite{MstarPaper}.
\textbf{N}, in both projects.
\end{itemize}
Also with W88 discharged: Theorem~S at $\vp$ (\lean{Unconditional.thmS_varpi}: \textbf{I}, \lean{Gen.NotLtilde}), Theorem~SE$_0$ at
$E_i$ (\lean{Bridge.thmSE0_varpi}: \textbf{U}), and the transcendence of $\pi/\vp^2$
(\lean{Bridge.transcendental_pi_div_varpi_sq}: \textbf{U}). Other CM periods and Beta values need their own period data in Lean.

\emph{What the labels rest on.} For \textbf{U} and \textbf{I} the group models are internal to the proof: the final statements
concern explicit numbers and the kernel checks the deduction, so they do not depend on the reviews of the models. One reads only
the definitions in the statements and, for \textbf{I}, accepts the inputs. For \textbf{W}, the reviews of \lean{lean_mstar} and of Section~\ref{sec:abvarv} apply. No human expert has reviewed any of this.

\emph{Provenance.} No new method is involved. Each entry is an instance of \cite[Thm.~4.1]{W88} or \cite[\S5d]{W88} on a group
$G''$ as above, and the verified entries go through \cite[\S5d]{W88}. For MG, SMG(a), SE(1), MX, MX$'$, deg2, X7, X9 and X$^\kappa$ the printed weight of Th\'eor\`eme~(iii) of
\cite{RoyWaldschmidt1993} already suffices on a product of standard groups (with W\"ustholz's theorem for the final statements),
and for MG$^\natural$ by an internal check in \cite{MstarPaper}. The proof in \cite{RoyWaldschmidt1993} reduces to $K=\QQ$ and
then refers to a refinement of an argument of Roy that is not written out (we could not consult Roy's paper). At $B/\pi$ and $B$
the printed weight gives no conclusion, and one needs $\alpha_{K,\Psi}$ of \cite{MstarPaper}. Theorem~M is classical in substance
\cite{MstarPaper}. For the other entries our claim is narrow: we found none of them stated in print, apart from what follows from
the conjugation arguments discussed in \cite{MstarPaper}, and we believe that the formal proofs here are the first of these
statements. The search was not exhaustive, this literature is moving fast and partly AI-generated, and \cite{MstarPaper} is itself
an unrefereed draft.

% ===========================================================================
\section{The artifact}\label{sec:stats}
% ===========================================================================

The formalization \cite{LSGRepo} uses Lean~4 \cite{Lean4} and Mathlib \cite{Mathlib}, version v4.35.0 (Mathlib commit
\lean{56a005855ea17e0b0958f73ec801ae140b0a1679}), with StrongSix \cite{StrongSixRepo} as a Lake git dependency pinned to commit \lean{b7f1649}. The default target
\lean{LSG} has 737 modules and about 168{,}500 lines. Table~\ref{tab:lines} gives the breakdown. The main theorem and the
vector-extension results of Section~\ref{sec:vecext} with everything they need take about 78{,}200 lines, with about
3{,}450 theorems and lemmas; the abelian chain alone (the main theorem and the abelian \S5d bridges) is 133 files and
about 35{,}300 lines. The abelian-surface layer of Section~\ref{sec:abvar} is 323 files and about 67{,}000 lines, of
which about 27{,}000 lines are generated certificate data. The vector extension of Section~\ref{sec:abvarv} adds 110
files and about 23{,}000 lines.

\begin{table}[ht]\footnotesize
\begin{tabular}{@{}lrp{0.47\textwidth}@{}}
\toprule
Directory & Lines & Content \\
\midrule
\multicolumn{3}{@{}l}{\emph{Abelian chain (default target)}}\\
\lean{Weierstrass/} & 3{,}768 & $\sigma$, $\zeta$, $\eta$, order two, Legendre, addition theorem, dictionary \\
\lean{SchneiderLang/}, \lean{Schneider*.lean} & 2{,}069 & Schneider--Lang criterion, Schneider's theorems, interfaces \\
\lean{LemniscateIndep/}, \lean{Bridge/} & 681 & $g_2(\vp\ZZ[i])=4$; statements for \lean{lean_mstar} \\
\lean{Subgroups/} & 399 & closed subgroups of $\RR^n$, subtori \\
\lean{Elliptic/AdditionLaws*} & 965 & addition laws on $E\subset\PP^2$ \\
\lean{ME/} & 16{,}516 & Stage 2: the multiplicity estimate \\
\lean{AuxFun/} (abelian) & 6{,}470 & Stage 3a: W83 Prop.~2.1, Lemma U \\
\lean{W88/} (abelian) & 6{,}971 & statement, Prop.~6.1 (750), \S7 (1{,}664), \S5d bridges and copies (4{,}004), sanity \\
\lean{CM/Coord}, \lean{CM/Transport} & 548 & Thm.~\ref{thm:W88} for abelian subvarieties \\
\midrule
\multicolumn{3}{@{}l}{\emph{Vector extensions (Section~\ref{sec:vecext})}}\\
\lean{MEV/} & 18{,}495 & multiplicity estimate for the class \\
\lean{AuxFunV/}, \lean{AuxFun/Prop21LatT*} & 5{,}320 & analytic side for $E^\natural$, polynomial fibre direction \\
\lean{Elliptic/Natural*} & 2{,}887 & the model of $E^\natural$ and its addition laws \\
\lean{W88/} (vector) & 12{,}465 & statements, Prop.~6.1, \S7, \S5d engine, bridges, copies \\
\lean{CM/} (rest) & 1{,}150 & CM toolkit for the bridges \\
\midrule
\multicolumn{3}{@{}l}{\emph{The Jacobian of $y^5=x(1-x)$ (Section~\ref{sec:abvar})}}\\
\lean{Jac/} & 36{,}518 & Grant's model, checkers and soundness (9{,}623); generated data and certificates (26{,}895) \\
\lean{JacAn/}, \lean{JacAn3/}, \lean{JacAn4/} & 9{,}256 & Abel map, Beta periods, exponential, $\theta$, laws on the image, surjectivity \\
\lean{MEB/} & 10{,}579 & multiplicity estimate for abelian blocks \\
\lean{JacBlock/} & 3{,}342 & the Jacobian block and its packages \\
\lean{AuxFunA/} & 1{,}777 & Prop.~2.1 under toric injectivity, Lemma U \\
\lean{W88/AbVar*} & 5{,}376 & statements, Prop.~6.1, \S7, \S5d engine, bridge, copies \\
\midrule
\multicolumn{3}{@{}l}{\emph{The vector extension $A_{\bar\sigma}^\natural$ (Section~\ref{sec:abvarv})}}\\
\lean{JacAnV/} & 2{,}244 & fibre coordinates, cocycle, fibre laws and periods \\
\lean{MEBV/} & 10{,}930 & multiplicity estimate for vector blocks \\
\lean{JacBlockV/} & 2{,}208 & the block $A_{\bar\sigma}^\natural$, packages, \S5d block properties \\
\lean{AuxFunAV/} & 1{,}976 & Prop.~2.1 under toric injectivity, Lemma U \\
\lean{W88/AbVarV*}, bridge & 5{,}115 & statements, Prop.~6.1, \S7, \S5d engine (4{,}297); bridge and copies (818) \\
\lean{Jac/B6V*} & 771 & Hilbert certificate (degree 41), Wronskian cocycle \\
\bottomrule
\end{tabular}
\caption{Lines per layer (\texttt{wc -l}), counted when each block was completed.}\label{tab:lines}
\end{table}

\emph{Verification.} The commands are
\begin{Verbatim}
lake exe cache get
lake build
lake env lean Axioms.lean
\end{Verbatim}
\lean{Axioms.lean} prints \texttt{\#print\ axioms} for 63 declarations, among them
\lean{LSG.W88.theorem41Abelian}, \lean{LSG.ME.multEstAb}, \lean{LSG.W88.prop61_b},
\lean{LSG.schneiderLang}, the Schneider theorems, \lean{LSG.Weierstrass.legendre}, \lean{LSG.CM.thm41Sub},
\lean{LSG.W88.thm11Kappa}, \lean{LSG.W88.hypStdCM}, the sanity derivations, and the vector-extension
endpoints of Section~\ref{sec:vecext} (\lean{LSG.MEV.multEstVec}, the five declarations named in Theorem~\ref{thm:vec}
and the six bridges), and those of Section~\ref{sec:abvar} (\lean{LSG.Jac.a1Facts}, \lean{LSG.MEB.multEstB_holds},
the two declarations named in Theorem~\ref{thm:abvar} and the two bridges), and those of Section~\ref{sec:abvarv}
(the three declarations named in Theorem~\ref{thm:abvarv} and the two bridges); every line reports \texttt{[propext,\ Classical.choice,\ Quot.sound]}. The source contains no \lean{sorry}, no
\lean{axiom} declaration and no \lean{native_decide}. A cold build of the default target, with Mathlib and StrongSix already built and \lean{LEAN_NUM_THREADS=4}, took
3\,h\,33\,min on a Windows machine with 32 logical processors and 31.7\,GB of RAM, with no errors. Times per file sum
to 14 hours, of which 9.4 are in \lean{LSG/Jac/}: the kernel certificates of Section~\ref{sec:abvar} dominate. Built two
at a time they take about five CPU-hours (about 2.7 hours wall-clock), with single files peaking at 4--5\,GB of memory;
at four at a time we sampled up to 7.6\,GB. An earlier attempt with Lake's default concurrency failed with
\lean{.olean} read errors, which we suspect but did not show to be memory exhaustion. Outside the
certificates, single modules took up to about nine minutes each.

Some vector-extension files contain polynomial identities whose cofactors were produced by Python
scripts (\lean{scripts/}); each identity is re-proved in Lean, so the scripts are untrusted inputs. These
files are in the default target; the abelian chain imports none of them and uses no generated data. The
certificates of Section~\ref{sec:abvar} are likewise generated by untrusted scripts and rechecked by the kernel.

% ===========================================================================
\section{Trust and limitations}\label{sec:trust}
% ===========================================================================

\emph{What the kernel guarantees.} That \lean{LSG.W88.theorem41Abelian} proves the proposition
\lean{Theorem41Abelian} from Lean's standard axioms. Whether that proposition is Waldschmidt's theorem
is a question about \lean{LSG/W88/Statement.lean}, reproduced in Section~\ref{sec:statement}, and has been
checked only by the AI review described there; likewise for the vector-extension statements and the review
of Section~\ref{sec:vecext}.

\emph{Scope.} The main theorem has \NR, and its $G_2$ is a product of elliptic curves (or, by
\lean{thm41Sub}, an abelian subvariety of one). The \S5d statement for distinct curves assumes
\lean{hrat}. For vector extensions only the cases of Theorem~\ref{thm:vec} are proved, and its premises
\lean{LiftAlg} and \lean{PerTanRat} are proved in Lean only for the CM curves used; the general
\lean{Theorem41Vec} is not proved. For the Jacobian of $y^5=x(1-x)$ only toric-injective hyperplanes and
period $A$-components are covered, and the identification of the formal torus with the Jacobian is
classical, not formalized (Section~\ref{sec:abvar}). For its vector extension the same two restrictions apply,
with $r=1$ and an adapted $b$; the identification of the model in $\PP^{11}$ with the extension is classical, not
formalized (Section~\ref{sec:abvarv}). The copies of \lean{lean_mstar}'s statements are checked by
typechecking inside \lean{lean_mstar}, which depends on this project; whether its group models describe
the intended groups is checked by its own reviews, not here.

\emph{Licence history.} Early in the project, about $1{,}500$ lines of the ancillary Lean files of
\cite{ChCh26} (the lemniscatic period identification) were ported for local use. Their licence did not
cover those files, so the port was never published; it was deleted and replaced by the independent proof
of Section~\ref{sec:stage1}, written by an agent that had not seen the ported files. The local history
still contains the port, and the repository will be published from a history that excludes it. Two
files of \lean{LSG/SchneiderLang/} adapt lemmas from StrongSix (same author, Apache~2.0); no other external
Lean code is included.

\emph{Reviews.} Apart from the Lean kernel, all checking so far was done by AI systems. No human
expert has read the statement against \cite{W88} or the proofs.

% ===========================================================================
\section{AI-assisted formalization}\label{sec:workflow}
% ===========================================================================

The Lean code, the planning documents and the text of this paper were produced by AI models (Claude,
Anthropic) working in Claude Code under the direction of the author, with several agent sessions in
parallel coordinated through written plans and interface files. The development git history, which is not published,
shows the main theorem and the vector-extension results taking about seventeen and a half hours of wall-clock time
(30--31 October 2026). Three practices from \cite{StrongSixPaper} were kept: every stage was fixed as a \lean{Prop}
interface first; each deduction between interfaces was proved before its inputs; and design problems
were recorded in the planning notes when found, together with the fix. The four problems of
Section~\ref{sec:stage2}, the irrelevant components in particular, were found by agents whose proofs failed
and were then traced to the mathematics, not to Lean. The author is responsible for the intended statements
and for the claims made here, and we would be grateful for corrections.

% ===========================================================================
\section*{Acknowledgements}
% ===========================================================================

We thank Michel Waldschmidt for making \cite{DALAG} and \cite{WaldschmidtLIL} freely available, the
Mathlib community for the library and in particular for its Weierstrass $\wp$ function and algebraic group
law on Weierstrass curves, and Nailin Guan and Yongle Hu, whose Cohen--Macaulay library enters through
StrongSix.

\bibliographystyle{amsplain}
\bibliography{paper}

\end{document}
